\chapter{ECS performance through cache and scheduling}\label{05-engine-i-ch017-ecs-performance-through-cache-and-scheduling}\phantomsection\hypertarget{05-engine-i-ch017}{}\label{stagechapter:05-engine-i-ch017}

\hypertarget{05-engine-i-ch017-ecs-performance-through-cache-and-scheduling}{}

\hyperref[volume-reading-map]{Reading map} · \hyperref[05-engine-i-ch016-archetypes-chunks-relationships-and-transforms]{Previous chapter} · \hyperref[05-engine-i-ch018-ecs-contracts-diagnostics-and-gpu-handoff]{Next chapter}

Data layout affects both memory traffic and parallel ownership, not just the number of component lookups. This chapter connects false sharing, synchronization, and diagnostics to the query patterns an ECS must serve.

\textbf{By the end:} Explain how cache-line ownership and scheduling influence an ECS workload, and propose measurements that isolate the bottleneck.

\textbf{Recall question:} How can two workers update different entities yet still contend on one cache line, and what layout or partition change would test that cause?

\textbf{False sharing in metadata}

Even if component data is perfectly partitioned, metadata can suffer.

Bad:

\begin{verbatim}
struct WorkerStats
{
    std::atomic<uint64_t> Count[8];
};
\end{verbatim}

If counters sit on one line, every worker invalidates the others' cache line.

Better:

\begin{verbatim}
struct alignas(64) WorkerCounter
{
    std::uint64_t Value = 0;
};
\end{verbatim}

then one per worker, followed by a reduction.

The exact destructive-interference size is architecture/platform dependent, but separating heavily written per-thread state is the principle.

\begin{center}\rule{0.5\linewidth}{0.5pt}\end{center}

\textbf{Read-only sharing is fine}

False sharing requires modifications.

Many threads reading the same cache line is generally not the same problem.

So:

\begin{verbatim}
immutable configuration
read-only Transform data
constant tables
\end{verbatim}

can be broadly shared.

The danger is write ownership bouncing.

\begin{center}\rule{0.5\linewidth}{0.5pt}\end{center}

\textbf{Prefetch and sequential iteration}

Modern CPUs are extremely good at recognizing sequential access.

Loop:

\begin{verbatim}
for (std::size_t i = 0; i < count; ++i)
{
    position[i] += velocity[i] * dt;
}
\end{verbatim}

has predictable streams:

\begin{verbatim}
position[i]
velocity[i]
\end{verbatim}

Hardware prefetchers can fetch future cache lines before they are demanded.

This is the payoff of dense arrays.

\begin{center}\rule{0.5\linewidth}{0.5pt}\end{center}

\textbf{Pointer chasing defeats prefetch}

Compare:

\begin{verbatim}
dense:
Component[0]
Component[1]
Component[2]
...
\end{verbatim}

with:

\begin{verbatim}
vector<Component*>
    ↓
random heap object
    ↓
random heap object
    ↓
random heap object
\end{verbatim}

The latter has poor spatial predictability.

Each pointer may stall waiting for another cache miss.

This is why ECS data layout can outperform a traditional object hierarchy even when algorithmic big-O is identical.

\begin{center}\rule{0.5\linewidth}{0.5pt}\end{center}

\textbf{Sparse-set query locality versus archetype locality}

Sparse-set join:

\begin{verbatim}
iterate Health dense array
      ↓
lookup Sparse<Transform>[entity]
      ↓
jump to Transform dense slot
      ↓
lookup Sparse<Velocity>[entity]
      ↓
jump to Velocity dense slot
\end{verbatim}

Each component pool is individually dense, but cross-component access order may be poorly correlated.

Archetype chunk:

\begin{verbatim}
Transform[row]
Velocity[row]
Health[row]
\end{verbatim}

all use the same row.

This produces excellent parallel sequential streams.

That is a core archetype advantage.

\begin{center}\rule{0.5\linewidth}{0.5pt}\end{center}

\textbf{Software prefetch}

You \emph{can} manually prefetch:

\begin{verbatim}
__builtin_prefetch(&components[i + distance]);
\end{verbatim}

on compilers that support it.

Do not reach for this first.

Sequential loops already benefit heavily from hardware prefetch.

Manual prefetch is mainly useful for irregular-but-predictable access patterns where profiling demonstrates memory stalls.

Poorly chosen prefetch can:

\begin{itemize}
\tightlist
\item
  waste bandwidth,
\item
  evict useful lines,
\item
  arrive too early,
\item
  arrive too late.
\end{itemize}

Data layout is more important than sprinkling prefetch intrinsics.

\begin{center}\rule{0.5\linewidth}{0.5pt}\end{center}

\textbf{System scheduling example for a real frame}

Consider these systems:

\begin{verbatim}
Input
CharacterController
AI
Animation
PhysicsPreSync
PhysicsStep
CollisionEvents
PhysicsPostSync
TransformHierarchy
RenderExtraction
AudioExtraction
\end{verbatim}

A possible graph:

\begin{verbatim}
Input ────────┐
              v
      CharacterController ──┐
                            v
AI --------------------> PhysicsPreSync
                            |
                            v
                       PhysicsStep
                            |
                  +---------+---------+
                  |                   |
                  v                   v
           CollisionEvents     PhysicsPostSync
                                      |
                                      v
                              TransformHierarchy
                              /               \
                             v                 v
                      RenderExtraction    AudioExtraction

Animation --------------------^
\end{verbatim}

Independent nodes can run in parallel when access declarations allow it.

This is where ECS stops being merely a storage library and becomes a \textbf{runtime scheduler architecture}.

\begin{center}\rule{0.5\linewidth}{0.5pt}\end{center}

\textbf{Read/write declaration example}

Imagine:

\begin{verbatim}
SystemDescriptor movement
{
    .Reads  = { Velocity },
    .Writes = { Transform }
};

SystemDescriptor renderer
{
    .Reads  =
    {
        Transform,
        Renderable
    }
};

SystemDescriptor ai
{
    .Reads  = { Transform },
    .Writes = { Velocity }
};
\end{verbatim}

Conflicts:

\begin{verbatim}
AI → Movement
\end{verbatim}

because AI writes Velocity that Movement reads.

And:

\begin{verbatim}
Movement → Renderer
\end{verbatim}

because Movement writes Transform that renderer reads.

AI and renderer both read Transform, but renderer also depends indirectly on movement.

The resulting DAG emerges from semantic and hazard edges together.

\begin{center}\rule{0.5\linewidth}{0.5pt}\end{center}

\textbf{Kairo's current Scene versus ECS runtime}

This distinction is crucial to your engine.

Current \BookInlineCode{{E}{n}{g}{i}{n}{e}{C}{o}{r}{e}{:}{:}{S}{c}{e}{n}{e}}{EngineCore::Scene} owns an authored record approximately like:

\begin{verbatim}
Record
├── Name
├── Transform
├── Settings
├── Parent
├── Children
├── optional MeshRenderer
├── optional SceneInstance
├── optional Camera
├── optional Light
├── optional AudioEmitter
├── optional AudioListener
├── optional Environment
├── optional Logic
├── optional RigidBody
└── optional Collider
\end{verbatim}

That is effectively an authored \textbf{entity record / AoS-like scene representation}.

It prioritizes:

\begin{itemize}
\tightlist
\item
  serialization,
\item
  editor semantics,
\item
  stable scene IDs,
\item
  explicit hierarchy,
\item
  validation.
\end{itemize}

That is not the same job as runtime ECS storage.

\begin{center}\rule{0.5\linewidth}{0.5pt}\end{center}

\textbf{Do not prematurely delete the authored Scene}

A strong architecture can be:

\begin{verbatim}
Authoring Scene
    |
    | extraction / bridge
    v
Runtime ECS
    |
    +------> PhysicsWorld
    |
    +------> RenderScene
    |
    +------> Audio
\end{verbatim}

The authored scene remains understandable and serializable.

The runtime ECS becomes optimized for simulation.

This is often superior to making the editor directly manipulate low-level chunk storage.

\begin{center}\rule{0.5\linewidth}{0.5pt}\end{center}

\textbf{Scene/ECS identity bridge}

Persistent scene identity:

\begin{verbatim}
Scene Entity ID
\end{verbatim}

Runtime identity:

\begin{verbatim}
ECS Entity(Index,Generation)
\end{verbatim}

Keep a bridge:

\begin{verbatim}
SceneID → ECS Entity
ECS Entity → SceneID
\end{verbatim}

Then:

\begin{itemize}
\tightlist
\item
  serialization uses SceneID,
\item
  simulation uses ECS Entity,
\item
  picking may use SceneID,
\item
  physics maps ECS Entity → BodyID,
\item
  rendering extracts from ECS.
\end{itemize}

Each subsystem gets identity suited to its lifetime.

\begin{center}\rule{0.5\linewidth}{0.5pt}\end{center}

\textbf{Current deferred runtime scene mutations}

Interestingly, \BookInlineCode{{E}{n}{g}{i}{n}{e}{C}{o}{r}{e}{:}{:}{R}{u}{n}{t}{i}{m}{e}{S}{c}{e}{n}{e}{C}{o}{m}{m}{a}{n}{d}{B}{u}{f}{f}{e}{r}}{EngineCore::RuntimeSceneCommandBuffer} already implements several ideas from Pass 10 for the authored Scene:

\begin{itemize}
\tightlist
\item
  pending entity tokens,
\item
  FIFO commands,
\item
  deferred create/destroy,
\item
  deferred parenting,
\item
  deferred tags,
\item
  explicit commit.
\end{itemize}

Its documentation says commands are applied only at synchronization points so callbacks/physics iteration do not invalidate scene traversal.

That is a very good architectural precedent.

But it is a \textbf{Scene command buffer}, not yet the \BookInlineCode{{K}{a}{i}{r}{o}{E}{C}{S}}{KairoECS} structural command buffer discussed in Pass 10.

The concept can be reused.

\begin{center}\rule{0.5\linewidth}{0.5pt}\end{center}

\textbf{Current runtime physics event design}

Your PhysicsWorld has:

\begin{verbatim}
enum class PhysicsContactEventType
{
    Begin,
    Stay,
    End
};
\end{verbatim}

and records:

\begin{verbatim}
BodyA
BodyB
ColliderA
ColliderB
Trigger
Response
Type
\end{verbatim}

This is already a strong subsystem-neutral event representation.

The next bridge should translate those runtime physics IDs into ECS entities.

Do not make gameplay search all bodies to discover their owner.

\subsection{Decide which pairs may interact, then publish overlap transitions safely}\label{05-engine-i-ch017-decide-which-pairs-may-interact-then-publish-overlap-transitions-safely}

\hypertarget{05-engine-i-ch017-decide-which-pairs-may-interact-then-publish-overlap-transitions-safely}{}

The existing \BookInlineCode{{P}{h}{y}{s}{i}{c}{s}{C}{o}{n}{t}{a}{c}{t}{E}{v}{e}{n}{t}{T}{y}{p}{e}}{PhysicsContactEventType} has \BookInlineCode{{B}{e}{g}{i}{n}}{Begin}, \BookInlineCode{{S}{t}{a}{y}}{Stay}, and \BookInlineCode{{E}{n}{d}}{End}, and the event record already carries body and collider identifiers plus trigger/response flags. The missing mechanism is how the engine decides that a pair is eligible, maintains its overlap state across ticks, and delivers a transition without exposing live physics storage to gameplay mutation. These are three separate decisions: \textbf{filtering} decides whether a pair may proceed; \textbf{narrow phase and simulation} decide what actually overlaps or touches; \textbf{event delivery} tells gameplay how the relation changed. A filter pass is not a collision test, and an overlap notification is not an impulse.

\subsubsection{A category and mask rule is reciprocal}\label{05-engine-i-ch017-a-category-and-mask-rule-is-reciprocal}

\hypertarget{05-engine-i-ch017-a-category-and-mask-rule-is-reciprocal}{}

Assign each shape a category bit set \(C\) and an accepted-category mask \(M\). For shapes \(A\) and \(B\), the ordinary rule is

\[
\operatorname{maskPass}(A,B) = (C_A \mathbin{\&} M_B \ne 0) \land (C_B \mathbin{\&} M_A \ne 0).
\]

Both directions must agree. If a player accepts walls but a wall rejects players, the pair is rejected by this policy. Categories describe what a shape \emph{is}; masks describe which categories that shape \emph{accepts}. Choose bits for interaction classes, not unique instances, and keep this rule at the broad-phase candidate boundary so rejected pairs do not reach narrow phase or solver construction. Other engine rules may also reject a pair, including body-type compatibility, same-body suppression, joints that disable connected-body collision, or a custom filter. State that policy separately; passing category and mask is not a promise that the pair will contact.

A small example uses \BookInlineCode{{P}{L}{A}{Y}{E}{R}{=}{0}{x}{1}}{PLAYER=0x1}, \BookInlineCode{{W}{O}{R}{L}{D}{=}{0}{x}{2}}{WORLD=0x2}, \BookInlineCode{{E}{N}{E}{M}{Y}{=}{0}{x}{4}}{ENEMY=0x4}, and \BookInlineCode{{P}{I}{C}{K}{U}{P}{=}{0}{x}{8}}{PICKUP=0x8}. Set the player's category to \BookInlineCode{{P}{L}{A}{Y}{E}{R}}{PLAYER} and mask to \BookInlineCode{{W}{O}{R}{L}{D}{\ }{\textbar{}}{\ }{E}{N}{E}{M}{Y}{\ }{\textbar{}}{\ }{P}{I}{C}{K}{U}{P}}{WORLD\ \textbar{}\ ENEMY\ \textbar{}\ PICKUP}. Set a pickup's category to \BookInlineCode{{P}{I}{C}{K}{U}{P}}{PICKUP} and mask to \BookInlineCode{{P}{L}{A}{Y}{E}{R}}{PLAYER}. Then both intersections are nonzero, so they pass. If the pickup mask becomes zero, the first intersection becomes zero while the second still passes; no pair is emitted. When a category is a bit set, one accepted shared bit suffices. A default mask must be deliberate: zero disables all ordinary pairs, while an all-bits mask accepts every category that the other shape permits.

\subsubsection{Same nonzero groups can override the masks}\label{05-engine-i-ch017-same-nonzero-groups-can-override-the-masks}

\hypertarget{05-engine-i-ch017-same-nonzero-groups-can-override-the-masks}{}

Box2D's documented rule gives a same, nonzero group index priority over the category/mask test: an equal positive group forces collision eligibility, and an equal negative group forbids it. If the group indexes differ or either is zero, evaluate the reciprocal mask rule. In pseudocode:

\begin{verbatim}
bool filter_allows(const Filter& a, const Filter& b) {
    // Box2D 3.1 policy: an equal nonzero group overrides both masks.
    if (a.group != 0 && a.group == b.group)
        return a.group > 0;

    // Each shape must accept the other shape's category.
    return (a.category & b.mask) != 0 &&
           (b.category & a.mask) != 0;
}
\end{verbatim}

For example, two parts assigned group \BookInlineCode{{+}{1}{2}}{+12} remain eligible even if both masks are zero; two parts assigned group \BookInlineCode{{-}{1}{2}}{-12} are rejected even if both masks accept everything. \BookInlineCode{{+}{1}{2}}{+12} and \BookInlineCode{{-}{1}{2}}{-12} are different groups, so their pair falls back to masks. This group rule is an engine-specific override convention, not a universal meaning of ``group.'' This example pins \textbf{Box2D 3.1.0} as a reference policy: its filter categories and masks are 64-bit, and it documents the equal nonzero group override above. Box2D 2.4 uses narrower category/mask fields, despite a similar group policy. Pin the engine/library version and adapt field widths and event APIs to that version; do not copy these semantics into another physics engine without checking its contract. The 3.1 policy also includes body-type rules beyond category/mask/group.

{\def\LTcaptype{none} % do not increment counter
\begin{longtable}[]{@{}
  >{\raggedright\arraybackslash}p{(\linewidth - 8\tabcolsep) * \real{0.1765}}
  >{\raggedleft\arraybackslash}p{(\linewidth - 8\tabcolsep) * \real{0.2353}}
  >{\raggedleft\arraybackslash}p{(\linewidth - 8\tabcolsep) * \real{0.2353}}
  >{\raggedright\arraybackslash}p{(\linewidth - 8\tabcolsep) * \real{0.1765}}
  >{\raggedright\arraybackslash}p{(\linewidth - 8\tabcolsep) * \real{0.1765}}@{}}
\toprule\noalign{}
\begin{minipage}[b]{\linewidth}\raggedright
Pair
\end{minipage} & \begin{minipage}[b]{\linewidth}\raggedleft
\(C_A \& M_B\)
\end{minipage} & \begin{minipage}[b]{\linewidth}\raggedleft
\(C_B \& M_A\)
\end{minipage} & \begin{minipage}[b]{\linewidth}\raggedright
Same nonzero group?
\end{minipage} & \begin{minipage}[b]{\linewidth}\raggedright
Result under the Box2D 3.1 filter rule
\end{minipage} \\
\midrule\noalign{}
\endhead
\bottomrule\noalign{}
\endlastfoot
Player / pickup & nonzero & nonzero & no & pass \\
Player / enemy with enemy mask zero & zero & nonzero & no & reject \\
Two \BookInlineCode{{+}{1}{2}}{+12} group parts, both masks zero & zero & zero & positive override & pass \\
Two \BookInlineCode{{-}{1}{2}}{-12} group parts, both masks all bits & nonzero & nonzero & negative override & reject \\
\BookInlineCode{{+}{1}{2}}{+12} part / \BookInlineCode{{-}{1}{2}}{-12} part with reciprocal bits & nonzero & nonzero & indexes differ & pass \\
\end{longtable}
}

A filter edit changes which pairs are eligible; it may invalidate contacts and generate end transitions. Treat it as a state change that needs reconciliation at a defined simulation boundary. Avoid changing filters in the middle of narrow-phase iteration or from a callback that is walking physics-owned arrays.

\subsubsection{A sensor detects overlap and skips physical response}\label{05-engine-i-ch017-a-sensor-detects-overlap-and-skips-physical-response}

\hypertarget{05-engine-i-ch017-a-sensor-detects-overlap-and-skips-physical-response}{}

A \textbf{sensor} uses the collision pipeline to ask whether shapes overlap, but it contributes no collision impulse or physical blocking response. A plate sensor can therefore register a crate without stopping it. A sensor still needs a geometry test and normally uses the same pair-filter policy; ``sensor'' does not mean ``ignore all masks.'' In the pinned v3.1.0 reference, sensor-event generation is disabled by default; enable it on the participating shape definitions according to that API. Its event path is conceptually:

\begin{verbatim}
conservative pair candidates
    -> category/mask/group and engine-rule filters
    -> exact overlap test
    -> record sensor overlap transition
    -> do not create a physical response row for this sensor pair
\end{verbatim}

The pinned Box2D 3.1.0 simulation documentation separates sensor events from contact events, disables sensor-event generation by default, and explicitly states sensors do not detect shapes that pass through them within one step. For a fast object or small sensor, it recommends a ray or shape cast. Do not generalize those CCD limits to later releases without checking their versioned contract. A trigger used only for gameplay should not be presented as a solid contact, and a physical contact may be relevant even when the game has no trigger behavior.

\subsubsection{Derive begin, stay, and end from stable pair sets}\label{05-engine-i-ch017-derive-begin-stay-and-end-from-stable-pair-sets}

\hypertarget{05-engine-i-ch017-derive-begin-stay-and-end-from-stable-pair-sets}{}

Store the set of logical overlaps from the last completed fixed step as \(P_k\) and the freshly validated overlap set for the new step as \(P_{k+1}\). Give each collider a stable identity that includes a slot or index and a generation, such as \BookInlineCode{{C}{o}{l}{l}{i}{d}{e}{r}{I}{d}{\{}{s}{l}{o}{t}{,}{\ }{g}{e}{n}{e}{r}{a}{t}{i}{o}{n}{\}}}{ColliderId\{slot,\ generation\}}. Form each unordered pair in canonical order, so \BookInlineCode{{(}{A}{,}{B}{)}}{(A,B)} and \BookInlineCode{{(}{B}{,}{A}{)}}{(B,A)} hash to the same key. Then:

\[
\begin{aligned}\mathrm{Begin}_{k+1} &= P_{k+1} \setminus P_k,\\\mathrm{Stay}_{k+1} &= P_{k+1} \cap P_k,\\\mathrm{End}_{k+1} &= P_k \setminus P_{k+1}.\end{aligned}
\]

These set equations avoid firing \BookInlineCode{{B}{e}{g}{i}{n}}{Begin} every frame an object remains in a trigger. The existing Kairo event enum calls the continuing state \BookInlineCode{{S}{t}{a}{y}}{Stay}; that is an application-level policy derived from set intersection. Box2D 3.1 documents buffered sensor/contact begin and end events (and contact hit events), not a universal stay callback. An engine may omit \BookInlineCode{{S}{t}{a}{y}}{Stay} to reduce event volume; if gameplay needs it, derive it from persistent state rather than assuming the physics library emits it.

A complete plate trace makes the transitions concrete. Let the plate be \BookInlineCode{{C}{o}{l}{l}{i}{d}{e}{r}{I}{d}{\{}{5}{,}{7}{\}}}{ColliderId\{5,7\}}, category \BookInlineCode{{P}{I}{C}{K}{U}{P}}{PICKUP}, mask \BookInlineCode{{P}{L}{A}{Y}{E}{R}}{PLAYER}; let the crate be \BookInlineCode{{C}{o}{l}{l}{i}{d}{e}{r}{I}{d}{\{}{2}{,}{4}{\}}}{ColliderId\{2,4\}}, category \BookInlineCode{{P}{L}{A}{Y}{E}{R}}{PLAYER}, mask \BookInlineCode{{P}{I}{C}{K}{U}{P}}{PICKUP}. Their reciprocal mask checks pass.

{\def\LTcaptype{none} % do not increment counter
\begin{longtable}[]{@{}
  >{\raggedleft\arraybackslash}p{(\linewidth - 6\tabcolsep) * \real{0.3077}}
  >{\raggedright\arraybackslash}p{(\linewidth - 6\tabcolsep) * \real{0.2308}}
  >{\raggedright\arraybackslash}p{(\linewidth - 6\tabcolsep) * \real{0.2308}}
  >{\raggedright\arraybackslash}p{(\linewidth - 6\tabcolsep) * \real{0.2308}}@{}}
\toprule\noalign{}
\begin{minipage}[b]{\linewidth}\raggedleft
Fixed tick
\end{minipage} & \begin{minipage}[b]{\linewidth}\raggedright
Validated overlap set
\end{minipage} & \begin{minipage}[b]{\linewidth}\raggedright
Derived event
\end{minipage} & \begin{minipage}[b]{\linewidth}\raggedright
Gameplay effect
\end{minipage} \\
\midrule\noalign{}
\endhead
\bottomrule\noalign{}
\endlastfoot
40 & empty & none & door remains closed \\
41 & \BookInlineCode{{\{}{(}{2}{:}{4}{,}{\ }{5}{:}{7}{)}{\}}}{\{(2:4,\ 5:7)\}} & \BookInlineCode{{B}{e}{g}{i}{n}}{Begin} & add crate 2:4 to plate occupancy; door begins opening \\
42 & \BookInlineCode{{\{}{(}{2}{:}{4}{,}{\ }{5}{:}{7}{)}{\}}}{\{(2:4,\ 5:7)\}} & \BookInlineCode{{S}{t}{a}{y}}{Stay} & keep occupancy; do not add a duplicate \\
43 & \BookInlineCode{{\{}{(}{2}{:}{4}{,}{\ }{5}{:}{7}{)}{\}}}{\{(2:4,\ 5:7)\}} & \BookInlineCode{{S}{t}{a}{y}}{Stay} & no repeated begin-side sound or reward \\
44 & empty because the crate leaves & \BookInlineCode{{E}{n}{d}{(}{S}{e}{p}{a}{r}{a}{t}{e}{d}{)}}{End(Separated)} & remove crate 2:4; close if the plate is now empty \\
\end{longtable}
}

Now replace tick 44 with a filter edit or destruction while overlapping. Reconciliation still emits an \BookInlineCode{{E}{n}{d}}{End}, with a reason such as \BookInlineCode{{F}{i}{l}{t}{e}{r}{C}{h}{a}{n}{g}{e}{d}}{FilterChanged} or \BookInlineCode{{D}{e}{s}{t}{r}{o}{y}{e}{d}}{Destroyed}. That reason lets the consumer distinguish ordinary motion from forced invalidation if its gameplay policy cares. Do not silently erase the pair from the saved set: that leaves a door count, pickup state, or audio loop latched. If a collider goes to sleep, sleep alone is not an overlap exit. Preserve its logical pair state; when it wakes or moves, revalidate geometry and filters, emitting a transition only if the relation actually changed. If an engine destroys contacts as part of sleep/wake or filter changes, normalize those implementation events back into logical pair state to avoid duplicate \BookInlineCode{{B}{e}{g}{i}{n}}{Begin} events.

\subsubsection{Queue copied events, then resolve at a safe boundary}\label{05-engine-i-ch017-queue-copied-events-then-resolve-at-a-safe-boundary}

\hypertarget{05-engine-i-ch017-queue-copied-events-then-resolve-at-a-safe-boundary}{}

Physics may run worker tasks or traverse internal arrays while producing contacts. Do not destroy a body, edit a collider, or mutate ECS structure from inside a narrow-phase, solver, or worker callback. During the physics step, collect compact event values (pair IDs, event kind, reason, and any small immutable measurement the consumer needs). Once the step and its storage traversal finish, map physics IDs to gameplay entities, dispatch read-only gameplay responses, and queue structural changes for an explicit synchronization point:

\begin{verbatim}
fixed-step input and forces
    -> physics step and pair reconciliation
    -> copy contact/sensor event values
    -> resolve generation-checked collider-to-entity mappings
    -> gameplay event systems
    -> deferred create/destroy/component/filter commands
    -> structural commit at the declared boundary
\end{verbatim}

The event queue owns its values. Never retain a pointer into a physics library's transient event arrays; Box2D documents its event arrays as invalidated by later world changes and warns that end events may refer to shapes or contacts already destroyed. If an end handler needs information after deletion, copy the stable gameplay owner ID and the small required payload into the event before processing the structural change. On resolution, check that the collider/entity generation still matches. If the old entity has already died, run cleanup using the copied old owner identity or discard according to explicit policy; never look up the slot alone and deliver the event to a new occupant.

A wake, deletion, teleport, body-type change, or filter edit can invalidate cached pair state. Reconcile at the fixed-step boundary and label forced end reasons when known. Deletion must emit or synthesize the end transition before removing the old pair from the logical set. A slot reused afterward increments its generation, so an old \BookInlineCode{{(}{s}{l}{o}{t}{=}{2}{,}{g}{e}{n}{e}{r}{a}{t}{i}{o}{n}{=}{4}{)}}{(slot=2,generation=4)} event cannot resolve to the new \BookInlineCode{{(}{s}{l}{o}{t}{=}{2}{,}{g}{e}{n}{e}{r}{a}{t}{i}{o}{n}{=}{5}{)}}{(slot=2,generation=5)} entity. This is the same identity rule used by generational registries: a memory slot is a location, not a lasting object identity.

\begin{allmightroleexercise}

\subsubsection{Exercise with a complete solution}\label{05-engine-i-ch017-exercise-with-a-complete-solution}

\hypertarget{05-engine-i-ch017-exercise-with-a-complete-solution}{}

A sensor gate has category \BookInlineCode{{G}{A}{T}{E}{=}{0}{x}{1}{0}}{GATE=0x10}, mask \BookInlineCode{{P}{L}{A}{Y}{E}{R}{\ }{\textbar{}}{\ }{K}{E}{Y}{\ }{=}{\ }{0}{x}{2}{1}}{PLAYER\ \textbar{}\ KEY\ =\ 0x21}, and group zero. A player has category \BookInlineCode{{P}{L}{A}{Y}{E}{R}{=}{0}{x}{1}}{PLAYER=0x1}, mask \BookInlineCode{{G}{A}{T}{E}{=}{0}{x}{1}{0}}{GATE=0x10}, group zero. A key has category \BookInlineCode{{K}{E}{Y}{=}{0}{x}{2}{0}}{KEY=0x20}, mask \BookInlineCode{{G}{A}{T}{E}{=}{0}{x}{1}{0}}{GATE=0x10}, group zero. A decoration has category \BookInlineCode{{D}{E}{C}{O}{R}{=}{0}{x}{4}{0}}{DECOR=0x40}, mask all bits, group zero. A pair of moving platform shapes both use group \BookInlineCode{{-}{5}}{-5}; a player/platform pair uses group indices \BookInlineCode{{-}{5}}{-5} and \BookInlineCode{{+}{5}}{+5}. Determine eligible pairs, then process the following gate observations: step 0 none; step 1 player overlaps gate; step 2 same overlap; step 3 player is destroyed; slot 3 is immediately reused for a decoration with a higher generation before queued gameplay events dispatch.

\end{allmightroleexercise}

\textbf{Solution.} Gate/player passes because \BookInlineCode{{G}{A}{T}{E}{\ }{\&}{\ }{p}{l}{a}{y}{e}{r}{.}{m}{a}{s}{k}{\ }{=}{\ }{0}{x}{1}{0}}{GATE\ \&\ player.mask\ =\ 0x10} and \BookInlineCode{{P}{L}{A}{Y}{E}{R}{\ }{\&}{\ }{g}{a}{t}{e}{.}{m}{a}{s}{k}{\ }{=}{\ }{0}{x}{1}}{PLAYER\ \&\ gate.mask\ =\ 0x1}; gate/key passes because \BookInlineCode{{G}{A}{T}{E}{\ }{\&}{\ }{k}{e}{y}{.}{m}{a}{s}{k}{\ }{=}{\ }{0}{x}{1}{0}}{GATE\ \&\ key.mask\ =\ 0x10} and \BookInlineCode{{K}{E}{Y}{\ }{\&}{\ }{g}{a}{t}{e}{.}{m}{a}{s}{k}{\ }{=}{\ }{0}{x}{2}{0}}{KEY\ \&\ gate.mask\ =\ 0x20}. Gate/decoration fails because \BookInlineCode{{D}{E}{C}{O}{R}{\ }{\&}{\ }{g}{a}{t}{e}{.}{m}{a}{s}{k}{\ }{=}{\ }{0}}{DECOR\ \&\ gate.mask\ =\ 0}, even though decoration accepts all categories. Two platform shapes with matching group \BookInlineCode{{-}{5}}{-5} are rejected regardless of masks. Player/platform groups \BookInlineCode{{-}{5}}{-5} and \BookInlineCode{{+}{5}}{+5} differ, so group override does not apply; use both reciprocal mask tests. If those bits are set, they pass this filter stage.

For the gate trace, step 0 has no event. Step 1 adds the stable pair \BookInlineCode{{(}{p}{l}{a}{y}{e}{r}{\ }{s}{l}{o}{t}{\ }{3}{,}{\ }{g}{e}{n}{e}{r}{a}{t}{i}{o}{n}{\ }{g}{;}{\ }{g}{a}{t}{e}{\ }{I}{D}{)}}{(player\ slot\ 3,\ generation\ g;\ gate\ ID)}, so emit exactly one \BookInlineCode{{B}{e}{g}{i}{n}}{Begin} and insert that generation-qualified player into the gate's occupancy set. Step 2 intersects the previous and current sets, so emit \BookInlineCode{{S}{t}{a}{y}}{Stay} only if this engine's gameplay layer requests it; do not emit another \BookInlineCode{{B}{e}{g}{i}{n}}{Begin}. At step 3 destruction removes the old pair and produces \BookInlineCode{{E}{n}{d}{(}{D}{e}{s}{t}{r}{o}{y}{e}{d}{)}}{End(Destroyed)} carrying the old generation and copied gameplay owner ID. The slot reuse creates a different identity \BookInlineCode{{(}{3}{,}{g}{+}{1}{)}}{(3,g+1)}. When the deferred queue drains, generation validation rejects the stale physics handle as a reference to the new decoration; cleanup uses the copied old owner data. The new decoration receives no gate event. The listener does not mutate the physics world while the event arrays or solver are being traversed.

\subsubsection{Bounded executable witness}\label{05-engine-i-ch017-bounded-executable-witness}

\hypertarget{05-engine-i-ch017-bounded-executable-witness}{}

The complete C++20 laboratory below implements this bounded filter policy, stable pair transitions, forced-end reasons and generation validation. Its checks remain active with \BookInlineCode{{-}{D}{N}{D}{E}{B}{U}{G}}{-DNDEBUG}; they also exercise a high category bit to detect unintended field narrowing. Compile from the vault root with \BookInlineCode{{c}{+}{+}{\ }{-}{s}{t}{d}{=}{c}{+}{+}{2}{0}{\ }{-}{O}{2}{\ }{-}{D}{N}{D}{E}{B}{U}{G}{\ }{-}{W}{a}{l}{l}{\ }{-}{W}{e}{x}{t}{r}{a}{\ }{-}{W}{p}{e}{d}{a}{n}{t}{i}{c}{\ }{-}{W}{e}{r}{r}{o}{r}{\ }{t}{o}{o}{l}{s}{/}{l}{a}{b}{s}{/}{c}{o}{l}{l}{i}{s}{i}{o}{n}{\_}{f}{i}{l}{t}{e}{r}{\_}{e}{v}{e}{n}{t}{\_}{l}{a}{b}{.}{c}{p}{p}{\ }{-}{o}{\ }{/}{t}{m}{p}{/}{c}{o}{l}{l}{i}{s}{i}{o}{n}{\_}{f}{i}{l}{t}{e}{r}{\_}{e}{v}{e}{n}{t}{\_}{l}{a}{b}}{c++\ -std=c++20\ -O2\ -DNDEBUG\ -Wall\ -Wextra\ -Wpedantic\ -Werror\ tools/labs/collision\_filter\_event\_lab.cpp\ -o\ /tmp/collision\_filter\_event\_lab}, then run \BookInlineCode{{/}{t}{m}{p}{/}{c}{o}{l}{l}{i}{s}{i}{o}{n}{\_}{f}{i}{l}{t}{e}{r}{\_}{e}{v}{e}{n}{t}{\_}{l}{a}{b}}{/tmp/collision\_filter\_event\_lab}. This is a deterministic policy model, not an actual Box2D/Kairo integration, a broad phase, a threaded dispatcher or a complete sleeping system. Implementation explanations are in the code comments.

\begin{verbatim}
// Bounded teaching witness for reciprocal category/mask filtering, same-group
// overrides, set-difference lifecycle events, and generation-safe event IDs.
// Compile: c++ -std=c++20 -O2 -Wall -Wextra -Wpedantic -Werror \
//   tools/labs/collision_filter_event_lab.cpp -o /tmp/collision_filter_event_lab
// Run: /tmp/collision_filter_event_lab
#include <algorithm>
#include <stdexcept>
#include <string>
#include <cstdint>
#include <iostream>
#include <map>
#include <set>
#include <tuple>
#include <vector>

// These checks remain active in release builds; NDEBUG must not erase the lab.
static unsigned checks = 0;
void check(bool condition) {
    ++checks;
    if (!condition) throw std::runtime_error("failed check " + std::to_string(checks));
}

struct Filter {
    std::uint64_t category;
    std::uint64_t mask;
    std::int32_t group;
};

// This models the Box2D 3.1 category/mask/group rule only.
// It does not include body-type, same-body, joint, or custom filtering rules.
bool should_collide(const Filter& a, const Filter& b) {
    if (a.group != 0 && a.group == b.group) {
        return a.group > 0; // Equal positive forces; equal negative forbids.
    }
    return (a.category & b.mask) != 0 && (b.category & a.mask) != 0;
}

struct ShapeId {
    std::uint32_t slot;
    std::uint32_t generation;
    friend bool operator==(ShapeId, ShapeId) = default;
    friend bool operator<(ShapeId a, ShapeId b) {
        return std::tie(a.slot, a.generation) < std::tie(b.slot, b.generation);
    }
};

struct Pair {
    ShapeId a;
    ShapeId b;
    static Pair canonical(ShapeId x, ShapeId y) {
        return y < x ? Pair{y, x} : Pair{x, y};
    }
    friend bool operator<(const Pair& x, const Pair& y) {
        return std::tie(x.a, x.b) < std::tie(y.a, y.b);
    }
};

enum class Kind { Begin, Stay, End };
enum class EndReason { Separated, FilterChanged, Destroyed, DiscontinuousMove };
struct Event {
    Kind kind;
    Pair pair;
    EndReason reason = EndReason::Separated;
};

// `previous` is logical touching state. Sleep does not erase it; a wake must
// revalidate geometry and filters before this function receives `current`.
std::vector<Event> transitions(const std::set<Pair>& previous,
                               const std::set<Pair>& current,
                               const std::map<Pair, EndReason>& forcedEnds = {}) {
    std::vector<Event> out;
    for (const Pair& p : current) {
        out.push_back({previous.contains(p) ? Kind::Stay : Kind::Begin, p});
    }
    for (const Pair& p : previous) {
        if (!current.contains(p)) {
            const auto it = forcedEnds.find(p);
            out.push_back({Kind::End, p, it == forcedEnds.end()
                                               ? EndReason::Separated
                                               : it->second});
        }
    }
    return out;
}

struct Slot {
    std::uint32_t generation = 1;
    bool alive = false;
};
bool valid(const std::vector<Slot>& slots, ShapeId id) {
    return id.slot < slots.size() && slots[id.slot].alive &&
           slots[id.slot].generation == id.generation;
}

int main() {
    constexpr std::uint64_t PLAYER = 1, ENEMY = 2, PLATE = 4, DEBRIS = 8;
    const Filter player{PLAYER, ENEMY | PLATE, 0};
    const Filter enemy{ENEMY, PLAYER, 0};
    const Filter plate{PLATE, PLAYER, 0};
    const Filter debris{DEBRIS, 0, 0};
    check(should_collide(player, enemy));        // both masks agree
    check(should_collide(player, plate));        // sensor-like category still filters
    check(!should_collide(player, debris));      // neither side accepts the other
    check(!should_collide(player, Filter{ENEMY, 0, 0})); // enemy rejects player
    check(should_collide(Filter{PLAYER, 0, 7}, Filter{ENEMY, 0, 7}));
    check(!should_collide(Filter{PLAYER, ~0ULL, -7}, Filter{ENEMY, ~0ULL, -7}));
    check(should_collide(Filter{PLAYER, ENEMY, 3}, Filter{ENEMY, PLAYER, -3}));
    check(should_collide(Filter{PLAYER, ENEMY, 3}, Filter{ENEMY, PLAYER, 0}));

    // A high category bit catches accidental narrowing to a 16/32-bit field.
    constexpr std::uint64_t HIGH = std::uint64_t{1} << 63;
    check(should_collide(Filter{HIGH, PLAYER, 0}, Filter{PLAYER, HIGH, 0}));
    check(transitions({}, {}).empty()); // no spurious event before any overlap

    const ShapeId oldActor{2, 4};
    const ShapeId sensor{5, 7};
    const Pair oldPair = Pair::canonical(oldActor, sensor);
    std::set<Pair> previous;
    std::set<Pair> current{oldPair};
    auto e = transitions(previous, current);
    check(e.size() == 1 && e[0].kind == Kind::Begin);
    previous = current;

    // Persisting overlap yields one application-level Stay; Box2D's buffered
    // event API documents begin/end and hit events, not a universal Stay event.
    e = transitions(previous, current);
    check(e.size() == 1 && e[0].kind == Kind::Stay);

    // A filter edit invalidates the old overlap and records why it ended.
    current.clear();
    e = transitions(previous, current, {{oldPair, EndReason::FilterChanged}});
    check(e.size() == 1 && e[0].kind == Kind::End &&
           e[0].reason == EndReason::FilterChanged);
    previous = current;

    // Re-enter, then destroy while overlapping. Copy the event IDs before any
    // structural mutation or callback; never keep engine-owned event pointers.
    current.insert(oldPair);
    e = transitions(previous, current);
    check(e.size() == 1 && e[0].kind == Kind::Begin);
    previous = current;
    current.clear();
    e = transitions(previous, current, {{oldPair, EndReason::Destroyed}});
    check(e.size() == 1 && e[0].pair.a == oldPair.a &&
           e[0].pair.b == oldPair.b && e[0].reason == EndReason::Destroyed);

    std::vector<Slot> slots(6);
    slots[2] = {4, true};
    slots[5] = {7, true};
    check(valid(slots, oldActor));
    // Deferred destruction runs at the safe structural boundary. Reusing slot
    // 2 changes generation, so an old event cannot target its new occupant.
    slots[2].alive = false;
    ++slots[2].generation;
    const ShapeId newActor{2, 5};
    slots[2].alive = true;
    check(!valid(slots, oldActor) && valid(slots, newActor));
    check(e[0].pair.a == oldPair.a && e[0].pair.a != newActor);

    std::cout << "collision filter/event lifecycle: " << checks << " checks PASS\n";
}
\end{verbatim}

\textbf{Versioned source policy.} The filter precedence and 64-bit widths above follow the official \href{https://raw.githubusercontent.com/erincatto/box2d/v3.1.0/docs/simulation.md}{Box2D 3.1.0 simulation documentation} and its \href{https://raw.githubusercontent.com/erincatto/box2d/v3.1.0/include/box2d/types.h}{tagged filter/type declarations}. The event-array lifetime and invalidation details follow its \href{https://raw.githubusercontent.com/erincatto/box2d/v3.1.0/include/box2d/box2d.h}{3.1.0 events reference}. These are a named reference policy for this example; a consuming project must pin and verify its actual engine version, because category widths, default event settings, same-body/body-type rules, sensor event enablement, callback constraints, and lifecycle APIs differ across releases and engines.

\begin{center}\rule{0.5\linewidth}{0.5pt}\end{center}

\textbf{Current runtime job design}

Your \BookInlineCode{{J}{o}{b}{S}{y}{s}{t}{e}{m}}{JobSystem} has worker threads and a queue but deliberately does not own:

\begin{verbatim}
renderer
physics
scene objects
\end{verbatim}

That is good.

Build ECS scheduling as:

\begin{verbatim}
SystemGraph
     ↓
ready systems
     ↓
JobSystem.Submit(...)
     ↓
barriers
     ↓
next dependency wave
\end{verbatim}

rather than stuffing ECS-specific logic into the generic thread pool.

\begin{center}\rule{0.5\linewidth}{0.5pt}\end{center}

\textbf{A production ECS scheduler descriptor}

A useful future design:

\begin{verbatim}
struct ECSSystemDescriptor
{
    SystemID ID;
    std::string Name;

    std::vector<ComponentTypeID> Reads;
    std::vector<ComponentTypeID> Writes;

    std::vector<ResourceTypeID> ResourceReads;
    std::vector<ResourceTypeID> ResourceWrites;

    std::vector<SystemID> Before;
    std::vector<SystemID> After;

    TickDomain Domain;

    bool MayQueueStructuralChanges = false;
};
\end{verbatim}

where:

\begin{verbatim}
enum class TickDomain
{
    Fixed,
    Variable,
    RenderExtraction
};
\end{verbatim}

Now the scheduler has enough information to reason about execution.

\begin{center}\rule{0.5\linewidth}{0.5pt}\end{center}

\textbf{Hazard function}

\begin{verbatim}
bool Intersects(
    std::span<const ComponentTypeID> a,
    std::span<const ComponentTypeID> b);

bool Conflicts(
    const ECSSystemDescriptor& a,
    const ECSSystemDescriptor& b)
{
    // A writes something B reads.
    if (Intersects(a.Writes, b.Reads))
        return true;

    // A writes something B writes.
    if (Intersects(a.Writes, b.Writes))
        return true;

    // B writes something A reads.
    if (Intersects(b.Writes, a.Reads))
        return true;

    // Repeat for world resources.
    if (Intersects(
            a.ResourceWrites,
            b.ResourceReads))
        return true;

    if (Intersects(
            a.ResourceWrites,
            b.ResourceWrites))
        return true;

    if (Intersects(
            b.ResourceWrites,
            a.ResourceReads))
        return true;

    return false;
}
\end{verbatim}

This is the core static hazard detector.

\begin{center}\rule{0.5\linewidth}{0.5pt}\end{center}

\textbf{But scheduler must not over-serialize everything}

Suppose two systems both write Transform but to provably disjoint entity groups:

\begin{verbatim}
PlayerTransformSystem
EnemyTransformSystem
\end{verbatim}

A component-type-only hazard model says conflict.

That is conservative but safe.

Later you can expose partition guarantees:

\begin{verbatim}
query filters
archetype ownership
entity ranges
exclusive tags
\end{verbatim}

to recover parallelism.

Start conservative.

Correct concurrency is far more valuable than theoretical maximal parallelism.

\begin{center}\rule{0.5\linewidth}{0.5pt}\end{center}

\textbf{Chunk change versions}

Archetype chunks enable a powerful optimization.

Suppose system writes Transform in chunk C.

Set:

\begin{verbatim}
C.TransformVersion = WorldVersion.
\end{verbatim}

Later renderer asks:

\begin{quote}
Did Transform change since my last update version?
\end{quote}

If not, skip work.

This is chunk-level change detection.

It cannot prove which individual entity changed, only:

\[
\text{something in this column/chunk may have changed}.
\]

But it is extremely cheap.

This mechanism can accelerate:

\begin{itemize}
\tightlist
\item
  transform extraction,
\item
  GPU uploads,
\item
  replication,
\item
  bounds updates,
\item
  renderer proxy updates.
\end{itemize}

\begin{center}\rule{0.5\linewidth}{0.5pt}\end{center}

\textbf{Previous/current transform and physics sync pipeline}

A robust fixed-tick flow:

\begin{verbatim}
Before fixed simulation:
PreviousTransform = CurrentTransform

Gameplay writes inputs/forces
        ↓
Physics Step
        ↓
physics body states
        ↓
PhysicsToTransform system
        ↓
CurrentTransform

Render:
alpha = accumulator / fixedDt
        ↓
interpolate Previous/Current
        ↓
RenderTransform
\end{verbatim}

Static/kinematic entities may follow different authority rules.

\begin{center}\rule{0.5\linewidth}{0.5pt}\end{center}

\textbf{Transform authority must be explicit}

Physics body modes:

\textbf{Static}

Scene/ECS Transform drives physics initialization.

No runtime physics integration.

\textbf{Kinematic}

Gameplay Transform generally drives physics.

Physics computes contacts but does not freely integrate body from forces.

\textbf{Dynamic}

Physics owns world motion.

After step:

\begin{verbatim}
Physics → Transform.
\end{verbatim}

If gameplay simultaneously writes Transform for a dynamic body, you have two authorities.

That must be defined.

Otherwise systems fight every frame.

\begin{center}\rule{0.5\linewidth}{0.5pt}\end{center}

\textbf{Render interpolation and physics authority}

Dynamic physics entity may therefore have:

\begin{verbatim}
PhysicsPreviousPose
PhysicsCurrentPose
RenderInterpolatedPose
\end{verbatim}

rather than overwriting authored local transform at arbitrary render cadence.

This avoids mixing simulation state and render state.

Again:

\[
\boxed{
\text{one logical object can legitimately have several state representations for different domains}.
}
\]

\begin{center}\rule{0.5\linewidth}{0.5pt}\end{center}

\textbf{Input tick timing}

Raw input is usually sampled once per rendered frame.

But fixed simulation may execute:

\begin{verbatim}
0 fixed ticks
1 fixed tick
2 fixed ticks
...
\end{verbatim}

during that frame.

``Pressed this frame'' events require care.

If JumpPressed is consumed by every fixed tick, two substeps could jump twice.

A robust input-command layer needs semantics like:

\begin{verbatim}
continuous:
Move axis
Look rate

edge-triggered:
JumpPressed
FirePressed
\end{verbatim}

with controlled consumption across fixed steps.

\begin{center}\rule{0.5\linewidth}{0.5pt}\end{center}

\textbf{Collision event timing}

Similarly:

\begin{verbatim}
Begin collision
\end{verbatim}

should normally be delivered once for a contact transition, not once per renderer frame.

PhysicsWorld already maintains Begin/Stay/End around simulation steps.

So collision event domain belongs naturally to fixed simulation time.

Rendering can observe resulting state later.

\begin{center}\rule{0.5\linewidth}{0.5pt}\end{center}

\textbf{Event queue lifetime}

Never put references into events that may outlive the producing storage.

Bad:

\begin{verbatim}
struct CollisionEvent
{
    const ContactManifold& Contact;
};
\end{verbatim}

if manifold vector is rebuilt next physics step.

Better:

\begin{verbatim}
struct CollisionEvent
{
    Entity A;
    Entity B;
    Vec3 Point;
    Vec3 Normal;
    float Impulse;
};
\end{verbatim}

or an immutable stable event record.

Events should own the data required by their lifetime.

\begin{center}\rule{0.5\linewidth}{0.5pt}\end{center}

\textbf{Render-picking lifetime}

Same principle.

Do not map GPU ObjectID to:

\begin{verbatim}
Transform*
\end{verbatim}

because component address may migrate before readback.

Map to stable identity:

\begin{verbatim}
SceneID
or
generational Entity
\end{verbatim}

and validate after readback.

Identity bridges repeatedly solve lifetime separation.

\begin{center}\rule{0.5\linewidth}{0.5pt}\end{center}

\textbf{False-sharing example with ECS counters}

Suppose every worker tracks:

\begin{verbatim}
struct Stats
{
    uint64_t Processed;
};

Stats stats[8];
\end{verbatim}

\BookInlineCode{{s}{i}{z}{e}{o}{f}{(}{S}{t}{a}{t}{s}{)}{=}{8}}{sizeof(Stats)=8}.

Eight counters may fit in one 64-byte cache line.

Eight cores increment different counters.

The line bounces constantly.

Instead:

\begin{verbatim}
struct alignas(64) WorkerStats
{
    uint64_t Processed = 0;
};
\end{verbatim}

then reduce afterward.

This is a textbook example where logical independence is not physical independence.

\begin{center}\rule{0.5\linewidth}{0.5pt}\end{center}

\textbf{False sharing with Transform columns}

Suppose Transform is 48 bytes.

Cache line 64 bytes.

Rows can straddle cache-line boundaries.

If job split happens at arbitrary row:

\begin{verbatim}
worker A ends at row 100
worker B starts row 101
\end{verbatim}

one or two boundary cache lines can be shared.

That's usually fine.

If workers alternate rows, almost every line is shared.

So \textbf{contiguous partitioning} matters much more than perfect cache-line-aligned boundaries.

\begin{center}\rule{0.5\linewidth}{0.5pt}\end{center}

\textbf{Prefetch and system fusion}

There is another performance consideration.

Suppose:

\begin{verbatim}
System A:
read Velocity
write Position

System B immediately after:
read Position
write Bounds
\end{verbatim}

Position written by A may still be hot in cache when B executes.

That can be good.

But combining every system into one giant loop destroys modularity and scheduling flexibility.

Sometimes system fusion improves locality.

Sometimes separation enables parallelism.

This is another benchmark decision, not a dogma.

\begin{center}\rule{0.5\linewidth}{0.5pt}\end{center}

\textbf{Archetype fragmentation}

Suppose entities differ by tiny structural flags:

\begin{verbatim}
A = Transform + Renderable
B = Transform + Renderable + Selected
C = Transform + Renderable + Highlighted
D = Transform + Renderable + Selected + Highlighted
\end{verbatim}

Now one logical population is split across four archetypes.

Queries must visit more chunks.

Small chunks waste capacity.

This is archetype fragmentation.

Watch metrics:

\[
\text{archetype count}
\]

\[
\text{chunk count}
\]

\[
\text{average chunk occupancy}
\]

\[
\text{wasted bytes}.
\]

Production ECS tools expose archetype/chunk occupancy for exactly this reason. Unity's Archetypes window, for example, reports chunk counts, chunk capacity and unused entity capacity per archetype. \href{https://docs.unity.cn/Packages/com.unity.entities\%401.0/manual/editor-archetypes-window.html?utm_source=chatgpt.com}{Unity User Manual}

\begin{center}\rule{0.5\linewidth}{0.5pt}\end{center}

\textbf{Chunk occupancy}

Define:

\[
Occupancy=
\frac{LiveEntities}
{TotalChunkCapacity}.
\]

Suppose archetype has ten chunks each capacity 256:

\[
capacity=2560.
\]

But only 900 entities:

\[
Occupancy=
900/2560
\approx35\%.
\]

Memory locality benefits are being squandered.

Structural churn/shared-value grouping/tag explosion can cause this.

\begin{center}\rule{0.5\linewidth}{0.5pt}\end{center}

\textbf{Sparse-set fragmentation is different}

Sparse sets do not have archetype fragmentation.

Each component type has one packed component array.

But cross-component query locality may be poorer.

So the trade is:

\begin{verbatim}
Sparse set:
excellent per-component density
less guaranteed multi-component alignment

Archetype:
excellent multi-component alignment
possible population fragmentation across signatures
\end{verbatim}

That is the actual nuanced comparison.

\begin{center}\rule{0.5\linewidth}{0.5pt}\end{center}

\textbf{Kairo recommendation for now}

Given current \BookInlineCode{{K}{a}{i}{r}{o}{E}{C}{S}}{KairoECS}, I would \textbf{not} replace the sparse-set core immediately.

Instead:

\begin{verbatim}
Stage 1
keep sparse sets
add generic queries
add command buffers
add scheduling metadata

Stage 2
benchmark real engine systems

Stage 3
if multi-component traversal dominates:
prototype archetype/chunk backend

Stage 4
compare objectively
\end{verbatim}

You need evidence from your engine workloads before paying archetype complexity.

\begin{center}\rule{0.5\linewidth}{0.5pt}\end{center}

\textbf{Good archetype prototype workload}

Build two storage implementations:

\textbf{Sparse set}

\begin{verbatim}
Transform pool
Velocity pool
Health pool
\end{verbatim}

\textbf{Archetype}

\begin{verbatim}
{Transform,Velocity,Health} chunks
\end{verbatim}

Benchmark:

\[
1,000,000
\]

entities.

Run:

\begin{verbatim}
integrate transforms
damage/health iteration
component add/remove churn
destroy/create churn
\end{verbatim}

Measure:

\begin{verbatim}
iteration ns/entity
structural mutation ns/entity
cache misses
bytes/entity
\end{verbatim}

Then you have data rather than architecture fashion.

\begin{center}\rule{0.5\linewidth}{0.5pt}\end{center}

\textbf{Repository audit --- M087/M088}

Current implementation status after inspecting KairoECS, KairoEngineCore, KairoPhysicsEngine and KairoRenderer:

{\def\LTcaptype{none} % do not increment counter
\begin{longtable}[]{@{}
  >{\raggedright\arraybackslash}p{(\linewidth - 2\tabcolsep) * \real{0.5000}}
  >{\raggedright\arraybackslash}p{(\linewidth - 2\tabcolsep) * \real{0.5000}}@{}}
\toprule\noalign{}
\begin{minipage}[b]{\linewidth}\raggedright
Concept
\end{minipage} & \begin{minipage}[b]{\linewidth}\raggedright
Current status
\end{minipage} \\
\midrule\noalign{}
\endhead
\bottomrule\noalign{}
\endlastfoot
Sparse-set ECS & Implemented \\
Archetype storage & Not currently in KairoECS \\
Archetype chunks & Not implemented \\
Archetype migration & Not implemented \\
AoS authored Scene records & Present in EngineCore Scene \\
SoA component pools & Present at KairoECS component level \\
Hot/cold runtime extraction & Partly present architecturally \\
Zero-size ECS tags & Not in KairoECS \\
Authored string tags & Implemented \\
World resources & Several systems naturally act as resources \\
Parent/child Scene hierarchy & Implemented \\
Cycle rejection & Implemented \\
World-transform hierarchy calculation & Implemented on demand \\
Dirty transform propagation & Not found in inspected ECS/Scene \\
Previous/current ECS transforms & Not found as a general transform system \\
Mesh renderer authored component & Implemented \\
Rigid body component & Implemented \\
Collider component & Implemented \\
Light component & Implemented \\
Camera component & Implemented \\
Logic/script authoring component & Implemented \\
Raw input state & Implemented as engine-level state \\
ECS input-intent component & Not generalised yet \\
Runtime physics IDs & Implemented in PhysicsWorld \\
Entity↔physics bridge & Architectural next step \\
Renderer ObjectID & Implemented \\
Physics Begin/Stay/End contacts & Implemented \\
Generic ECS event queue & Not yet a dedicated KairoECS layer \\
Deferred Scene mutations & Implemented in EngineCore \\
Fixed physics stepping & Implemented \\
Generic thread pool & Implemented \\
ECS dependency DAG & Not currently implemented \\
ECS access declarations & Not implemented \\
Automatic hazard detection & Not implemented \\
ECS topological scheduler & Not implemented \\
Parallel chunk iteration & Not applicable until chunk backend exists \\
\end{longtable}
}

\begin{center}\rule{0.5\linewidth}{0.5pt}\end{center}

\textbf{Reconstruction challenge --- archetype ECS}

From memory, derive this chain:

\begin{verbatim}
Entity has signature S
      ↓
lookup Archetype(S)
      ↓
Archetype owns chunks
      ↓
chunk owns component columns
      ↓
entity location = archetype + chunk + row
\end{verbatim}

Adding component C:

\begin{verbatim}
S
 ↓
S ∪ {C}
 ↓
destination archetype
 ↓
allocate destination row
 ↓
move common columns
 ↓
construct C
 ↓
update entity location
 ↓
swap-remove source row
 ↓
repair moved source entity location
\end{verbatim}

If you cannot explain all of those steps and why each is needed, archetype ECS is not yet mastered.

\begin{center}\rule{0.5\linewidth}{0.5pt}\end{center}

\textbf{Reconstruction challenge --- sparse set versus archetype}

You should be able to defend both positions.

Sparse set:

\[
\text{cheap structural component changes}
+
\text{excellent independent component density}
+
\text{O(1) lookup}
\]

but:

\[
\text{multi-component queries require membership lookups}.
\]

Archetype:

\[
\text{excellent multi-component iteration}
+
\text{query matching at archetype granularity}
+
\text{natural chunk jobs}
\]

but:

\[
\text{component add/remove = entity migration}.
\]

There is no correct answer without workload.

\begin{center}\rule{0.5\linewidth}{0.5pt}\end{center}

\textbf{Reconstruction challenge --- transform hierarchy}

Given:

\begin{verbatim}
A root
└── B
    └── C
\end{verbatim}

derive:

\[
W_A=L_A
\]

\[
W_B=W_AL_B
\]

\[
W_C=W_BL_C.
\]

Then explain why parent must execute before child.

Then explain two dirty models:

\begin{verbatim}
eager subtree dirty marking
\end{verbatim}

versus:

\begin{verbatim}
lazy parent-version comparison.
\end{verbatim}

Then explain fixed-tick interpolation:

\[
T_r=
Interpolate(
T_{prev},
T_{curr},
\alpha).
\]

That completes M087 conceptually.

\begin{center}\rule{0.5\linewidth}{0.5pt}\end{center}

\textbf{Reconstruction challenge --- engine subsystem bridge}

You should be able to trace one entity through the whole engine:

\begin{verbatim}
SceneID 57
    ↓
runtime ECS Entity {231,4}
    ↓
RigidBodyComponent
ColliderComponent
Transform
Renderable
    ↓
physics bridge:
Entity {231,4} ↔ BodyID 81
    ↓
PhysicsContactEvent Body81/Body99
    ↓
BodyID→Entity translation
    ↓
ECS CollisionEvent
    ↓
gameplay
\end{verbatim}

Simultaneously:

\begin{verbatim}
Entity {231,4}
    ↓
Renderable + Transform
    ↓
render extraction
    ↓
MeshDraw
ObjectID
    ↓
GPU picking
    ↓
ID→entity lookup
\end{verbatim}

If those identity domains are kept separate, the engine remains robust.

\begin{center}\rule{0.5\linewidth}{0.5pt}\end{center}

\textbf{Reconstruction challenge --- scheduler}

Given:

\begin{verbatim}
A reads Position, writes Velocity
B reads Velocity, writes Transform
C reads Transform
D reads Position
\end{verbatim}

You should immediately identify:

\[
A\rightarrow B
\]

because A writes what B reads,

and:

\[
B\rightarrow C
\]

because B writes what C reads.

A and D can potentially run concurrently if no hidden resource dependencies exist.

Then topologically schedule:

\begin{verbatim}
Wave 0: A, D
Wave 1: B
Wave 2: C
\end{verbatim}

That is M088's central scheduling skill.

\begin{center}\rule{0.5\linewidth}{0.5pt}\end{center}

\textbf{Reconstruction challenge --- false sharing}

Explain why this can be race-free but slow:

\begin{verbatim}
core 0 writes values[0]
core 1 writes values[1]
\end{verbatim}

when both occupy one cache line.

Then explain why:

\begin{verbatim}
core 0 writes [0..999]
core 1 writes [1000..1999]
\end{verbatim}

usually behaves far better.

If you understand the difference between:

\[
\text{data race}
\]

and:

\[
\text{cache-coherence contention},
\]

you understand false sharing.

\begin{center}\rule{0.5\linewidth}{0.5pt}\end{center}

\textbf{Exact coverage}

{\def\LTcaptype{none} % do not increment counter
\begin{longtable}[]{@{}rll@{}}
\toprule\noalign{}
Part & Topic & Covered \\
\midrule\noalign{}
\endhead
\bottomrule\noalign{}
\endlastfoot
1435 & Archetype ECS concept & ✓ \\
1436 & Archetype signatures & ✓ \\
1437 & Archetype chunks & ✓ \\
1438 & Entity movement between archetypes & ✓ \\
1439 & Chunk capacity and alignment & ✓ \\
1440 & Columnar component storage & ✓ \\
1441 & Sparse-set versus archetype trade-offs & ✓ \\
1442 & AoS component data & ✓ \\
1443 & SoA component data & ✓ \\
1444 & Hot-and-cold component splitting & ✓ \\
1445 & Tag components & ✓ \\
1446 & Singleton/resource components & ✓ \\
1447 & Relationship components & ✓ \\
1448 & Parent-child relationships & ✓ \\
1449 & Scene hierarchy transforms & ✓ \\
1450 & Dirty transform propagation & ✓ \\
1451 & Previous/current transform snapshots & ✓ \\
1452 & Renderable component & ✓ \\
1453 & Physics body component & ✓ \\
1454 & Collider component & ✓ \\
1455 & Light component & ✓ \\
1456 & Camera component & ✓ \\
1457 & Script/behavior component & ✓ \\
1458 & Input component/state & ✓ \\
1459 & Entity-to-physics mapping & ✓ \\
1460 & Entity-to-render-object mapping & ✓ \\
1461 & Stable picking IDs & ✓ \\
1462 & ECS event queues & ✓ \\
1463 & Collision-event delivery & ✓ \\
1464 & Begin-play lifecycle & ✓ \\
1465 & Fixed tick vs render tick & ✓ \\
1466 & System ordering dependencies & ✓ \\
1467 & System dependency graphs & ✓ \\
1468 & Topological scheduling & ✓ \\
1469 & Parallel-safe ECS iteration & ✓ \\
1470 & Read-only vs writable declarations & ✓ \\
1471 & Data-hazard detection & ✓ \\
1472 & ECS job partitioning & ✓ \\
1473 & False sharing & ✓ \\
1474 & Prefetch/sequential iteration & ✓ \\
\end{longtable}
}

The architecture you should now hold mentally is no longer merely:

\[
Entity\rightarrow Component.
\]

It is:

\[
\boxed{
\begin{aligned}
&\text{persistent scene identity}\\
&\qquad\downarrow\\
&\text{generational runtime entity}\\
&\qquad\downarrow\\
&\text{sparse pool or archetype/chunk storage}\\
&\qquad\downarrow\\
&\text{system queries}\\
&\qquad\downarrow\\
&\text{read/write declarations}\\
&\qquad\downarrow\\
&\text{dependency DAG}\\
&\qquad\downarrow\\
&\text{parallel jobs}\\
&\qquad\downarrow\\
&\text{renderer / physics / audio bridges}.
\end{aligned}
}
\]

And throughout the whole chain, \textbf{identity, storage location, subsystem handles and persistent IDs are deliberately separate things}. That separation is what lets data move freely without destroying correctness.

\textbf{Next: Pass 12 --- M089 + M090, Parts 1475--1500.} That finishes the ECS phase with the ugly production details that toy ECS tutorials skip: \textbf{non-trivial C++ component lifetime, move-only components, serialization, entity-reference remapping, schema migration, reflection/editor inspection, undo/redo, prefab data, scene→ECS loading, cloning, network/determinism concerns, ECS diagnostics, dangling-pointer detection, million-entity memory budgets, complexity benchmarking and finally implementing/defending the entire generational sparse-set ECS from a blank file.}

\subsection{Component Lifetime, Move-Only Data, Serialization, Entity Reference Remapping, Schema Evolution, Reflection, Inspector Editing, Undo/Redo, Prefabs, Scene↔ECS Conversion, Cloning, Network Identity, Determinism, ECS Diagnostics, Memory Budgets, Complexity, Invalidation Contracts, Reuse Testing, Benchmarking, and Building the Whole Sparse-Set ECS From Scratch}\label{05-engine-i-ch017-component-lifetime-move-only-data-serialization-entity-reference-remapping-schema-evolution-reflection-inspector-editing-undoredo-prefabs-sceneecs-conversion-cloning-network-identity-determinism-ecs-diagnostics-memory-budgets-complexity-invalidation-contracts-reuse-testing-benchmarking-and-building-the-whole-sparse-set-ecs-from-scratch}

\hypertarget{05-engine-i-ch017-component-lifetime-move-only-data-serialization-entity-reference-remapping-schema-evolution-reflection-inspector-editing-undoredo-prefabs-sceneecs-conversion-cloning-network-identity-determinism-ecs-diagnostics-memory-budgets-complexity-invalidation-contracts-reuse-testing-benchmarking-and-building-the-whole-sparse-set-ecs-from-scratch}{}

This pass \textbf{finishes Phase 15}.

Passes 10 and 11 gave us the structural model:

\[
Entity
\rightarrow
ComponentStorage
\rightarrow
Systems.
\]

Pass 12 asks what happens when those components are \textbf{real C++ objects}, the world must survive save/load/editor/network boundaries, and somebody asks you:

\begin{quote}
``Why did you choose this exact data structure, what invalidates what, what is the memory cost at one million entities, and can you prove stale handles never alias?''
\end{quote}

That is where an ECS stops being a coding exercise.

\begin{center}\rule{0.5\linewidth}{0.5pt}\end{center}

\subsection{PRODUCTION COMPONENT LIFETIME, SERIALIZATION, REFLECTION, SCENES AND IDENTITY}\label{05-engine-i-ch017-production-component-lifetime-serialization-reflection-scenes-and-identity}

\hypertarget{05-engine-i-ch017-production-component-lifetime-serialization-reflection-scenes-and-identity}{}

\begin{center}\rule{0.5\linewidth}{0.5pt}\end{center}

\textbf{Component construction and destruction}

So far we have often used components like:

\begin{verbatim}
struct Position
{
    float x;
    float y;
    float z;
};
\end{verbatim}

For such a type, storage looks deceptively simple.

Allocate bytes.

Write three floats.

Eventually release bytes.

But C++ components are not just byte ranges.

Every object has a \textbf{lifetime}.

For an object \(T\), there is a point at which suitably aligned storage exists but \textbf{no T object exists yet}, then construction begins its lifetime, and eventually destruction ends that lifetime. C++ object lifetime and storage duration are related but not identical. \href{https://en.cppreference.com/cpp/language/lifetime?utm_source=chatgpt.com}{Cppreference}

This distinction becomes critical the moment an ECS owns raw chunk storage.

\begin{center}\rule{0.5\linewidth}{0.5pt}\end{center}

\textbf{Storage is not an object}

Suppose an archetype chunk allocates:

\begin{verbatim}
std::byte memory[16384];
\end{verbatim}

and reserves 64 bytes at some offset for a component.

Those 64 bytes do \textbf{not} magically contain a valid:

\begin{verbatim}
std::string
\end{verbatim}

just because:

\begin{verbatim}
sizeof(std::string) <= 64.
\end{verbatim}

You must begin the object's lifetime.

For ordinary dynamic construction:

\begin{verbatim}
std::construct_at(
    reinterpret_cast<T*>(address),
    constructorArguments...);
\end{verbatim}

Later:

\begin{verbatim}
std::destroy_at(
    reinterpret_cast<T*>(address));
\end{verbatim}

The sequence is:

\[
\boxed{
Storage
\rightarrow
Construction
\rightarrow
Live\ T
\rightarrow
Destruction
\rightarrow
Storage.
}
\]

You cannot legally treat arbitrary bytes as a fully initialized nontrivial T.

\begin{center}\rule{0.5\linewidth}{0.5pt}\end{center}

\textbf{Sparse sets hide much of this for you}

Your current Kairo sparse set stores:

\begin{verbatim}
std::vector<Component> m_Components;
\end{verbatim}

That is important.

\BookInlineCode{{s}{t}{d}{:}{:}{v}{e}{c}{t}{o}{r}{\textless{}}{T}{\textgreater{}}}{std::vector\textless{}T\textgreater{}} already manages:

\begin{itemize}
\tightlist
\item
  allocation,
\item
  alignment,
\item
  construction,
\item
  move/copy during reallocation,
\item
  destruction,
\item
  exception cleanup.
\end{itemize}

So:

\begin{verbatim}
m_Components.emplace_back(...);
\end{verbatim}

constructs a real Component.

And:

\begin{verbatim}
m_Components.pop_back();
\end{verbatim}

destroys one.

This is a major simplification compared with writing a raw archetype chunk allocator.

\begin{center}\rule{0.5\linewidth}{0.5pt}\end{center}

\textbf{Archetype chunks force you to become a miniature \BookInlineCode{{s}{t}{d}{:}{:}{v}{e}{c}{t}{o}{r}}{std::vector}}

If you later implement raw archetype chunks, \textbf{you inherit responsibility for object lifetime}.

An archetype column cannot merely know:

\begin{verbatim}
component size = 48
alignment = 16
\end{verbatim}

It also needs operations.

Conceptually:

\begin{verbatim}
struct ComponentOperations
{
    std::size_t Size;
    std::size_t Alignment;

    void (*DefaultConstruct)(void*);
    void (*MoveConstruct)(void*, void*);
    void (*Destroy)(void*);

    // Depending on policy:
    void (*CopyConstruct)(void*, const void*);
};
\end{verbatim}

Then moving an entity from one archetype to another becomes type-erased lifetime management.

\begin{center}\rule{0.5\linewidth}{0.5pt}\end{center}

\textbf{Type-erased construction}

For type T:

\begin{verbatim}
template<typename T>
ComponentOperations MakeOperations()
{
    return
    {
        .Size = sizeof(T),
        .Alignment = alignof(T),

        .DefaultConstruct =
            [](void* destination)
            {
                std::construct_at(
                    static_cast<T*>(destination));
            },

        .MoveConstruct =
            [](void* destination,
               void* source)
            {
                std::construct_at(
                    static_cast<T*>(destination),
                    std::move(
                        *static_cast<T*>(source)));
            },

        .Destroy =
            [](void* object)
            {
                std::destroy_at(
                    static_cast<T*>(object));
            }
    };
}
\end{verbatim}

Now your ECS can manipulate unknown component types through their descriptors.

This is the runtime type-erasure layer archetype storage ultimately requires.

\begin{center}\rule{0.5\linewidth}{0.5pt}\end{center}

\textbf{Destruction order matters}

Imagine entity owns:

\begin{verbatim}
Transform
ScriptRuntime
AudioSource
PhysicsProxy
\end{verbatim}

If some component destructors refer to other runtime state, arbitrary destruction ordering can become observable.

Ideally components should not depend heavily on one another in destructors.

Instead subsystem teardown should occur in explicit systems:

\begin{verbatim}
pre-destroy phase
    ↓
release physics proxy
release script instance
release renderer proxy
    ↓
destroy ECS components
\end{verbatim}

This prevents hidden destructor ordering dependencies.

That is another data-oriented principle:

\[
\boxed{
\text{behavioral lifecycle should usually be explicit,
not buried in destructors}.
}
\]

\begin{center}\rule{0.5\linewidth}{0.5pt}\end{center}

\textbf{Component destruction on entity destruction}

For sparse-set Kairo:

\begin{verbatim}
Destroy Entity
    ↓
for every storage:
    Remove(entity)
    ↓
component swap-remove
    ↓
component destructor eventually runs on removed tail object
    ↓
generation increments
    ↓
slot recycled
\end{verbatim}

This ordering from Pass 10 remains important.

Runtime identity cannot be reused until old component objects have ceased to exist logically.

\begin{center}\rule{0.5\linewidth}{0.5pt}\end{center}

\textbf{Non-trivial C++ component types}

A trivial-looking component:

\begin{verbatim}
struct Health
{
    float Value;
};
\end{verbatim}

does not own resources.

A non-trivial component may contain:

\begin{verbatim}
struct Inventory
{
    std::vector<Item> Items;
    std::string OwnerName;
};
\end{verbatim}

or:

\begin{verbatim}
struct ScriptRuntime
{
    std::shared_ptr<VM> Machine;
    std::vector<Value> Locals;
};
\end{verbatim}

or:

\begin{verbatim}
struct AudioEmitterRuntime
{
    AudioVoiceHandle Voice;
    std::string DebugName;
};
\end{verbatim}

These have real construction/destruction semantics.

\begin{center}\rule{0.5\linewidth}{0.5pt}\end{center}

\textbf{Why raw \BookInlineCode{{m}{e}{m}{c}{p}{y}}{memcpy} becomes illegal or semantically wrong}

Suppose:

\begin{verbatim}
std::string a = "Kairo";
\end{verbatim}

and you byte-copy its object representation into some other location:

\begin{verbatim}
std::memcpy(&b, &a, sizeof(std::string));
\end{verbatim}

Now two string objects may contain internal pointers referring to the same allocation without correct ownership semantics.

Destruction can double-free.

The broader rule:

\[
\boxed{
\text{object relocation is not equivalent to byte copying
unless the type's semantics permit it}.
}
\]

Trivially copyable types can be handled much more aggressively.

Nontrivial types require constructors/move operations/destructors.

\begin{center}\rule{0.5\linewidth}{0.5pt}\end{center}

\textbf{Traits should affect storage policy}

A sophisticated ECS can branch:

\begin{verbatim}
if constexpr (
    std::is_trivially_copyable_v<T>)
{
    // memcpy/memmove may be available
}
else
{
    // actual move construction/destruction
}
\end{verbatim}

Likewise:

\begin{verbatim}
std::is_trivially_destructible_v<T>
\end{verbatim}

can eliminate destructor loops.

For a million POD-like physics components, skipping unnecessary lifetime callbacks can materially matter.

\begin{center}\rule{0.5\linewidth}{0.5pt}\end{center}

\textbf{Trivial relocation is broader than trivial copying conceptually}

Engines frequently want:

\begin{quote}
``This type can safely be relocated from one storage address to another.''
\end{quote}

That is not exactly the same question as:

\begin{verbatim}
std::is_trivially_copyable_v<T>.
\end{verbatim}

For example a type could own a heap allocation through a pointer and moving its raw bytes might be semantically acceptable if the old object's lifetime is appropriately ended---but standard C++ has historically made general relocatability subtler than engine programmers might wish.

Unless you establish a stronger engine-specific relocation contract, use actual move construction for nontrivial types.

\begin{center}\rule{0.5\linewidth}{0.5pt}\end{center}

\textbf{Your current Kairo removal contract has an important consequence}

Current swap-removal essentially performs:

\begin{verbatim}
m_Components[dense] =
    std::move(m_Components[last]);
\end{verbatim}

Therefore Component must support:

\[
\text{move assignment}.
\]

More precisely, because the current \BookInlineCode{{R}{e}{m}{o}{v}{e}{(}{)}}{Remove()} path is \BookInlineCode{{n}{o}{e}{x}{c}{e}{p}{t}}{noexcept}, a throwing move assignment can lead to:

\[
std::terminate.
\]

So the practical contract should be stated as something close to:

\begin{verbatim}
static_assert(
    std::is_nothrow_move_assignable_v<Component>);
\end{verbatim}

or Kairo should weaken its \BookInlineCode{{n}{o}{e}{x}{c}{e}{p}{t}}{noexcept} guarantee.

This is not academic.

It is an API contract that users of the ECS deserve to know.

\begin{center}\rule{0.5\linewidth}{0.5pt}\end{center}

\textbf{A stronger sparse-set implementation can move-construct instead}

One alternative storage model is raw/paged storage.

To remove row \(d\):

\begin{verbatim}
destroy object at d
move-construct object from last into d
destroy moved-from last
\end{verbatim}

rather than move-assigning into an existing object.

Conceptually:

\begin{verbatim}
std::destroy_at(destination);

std::construct_at(
    destination,
    std::move(*last));

std::destroy_at(last);
\end{verbatim}

But exception safety becomes harder.

If move construction throws after destination destruction, you have a hole.

So once again:

\[
\boxed{
nothrow move construction
}
\]

is extremely valuable for ECS components.

\begin{center}\rule{0.5\linewidth}{0.5pt}\end{center}

\textbf{Nontrivial components can hide allocations}

Consider:

\begin{verbatim}
struct Name
{
    std::string Value;
};
\end{verbatim}

Dense ECS iteration over one million Names is technically contiguous in terms of \BookInlineCode{{s}{t}{d}{:}{:}{s}{t}{r}{i}{n}{g}}{std::string} objects.

But each long string may point somewhere else in the heap.

So:

\begin{verbatim}
dense component array
\end{verbatim}

does not automatically imply:

\begin{verbatim}
all logically relevant data is contiguous.
\end{verbatim}

Hot components should avoid unnecessary dynamic allocations.

Cold authoring metadata can tolerate them much more readily.

\begin{center}\rule{0.5\linewidth}{0.5pt}\end{center}

\textbf{Move-only components}

Some engine objects should have exactly one owner.

Example:

\begin{verbatim}
struct ScriptInstance
{
    std::unique_ptr<VirtualMachine> VM;
};
\end{verbatim}

or:

\begin{verbatim}
struct GPUUploadTicket
{
    std::unique_ptr<UploadPayload> Payload;
};
\end{verbatim}

These are move-only.

\begin{verbatim}
ScriptInstance(const ScriptInstance&) = delete;
ScriptInstance& operator=(
    const ScriptInstance&) = delete;
\end{verbatim}

but move is allowed.

An ECS should not unnecessarily require components to be copyable.

\begin{center}\rule{0.5\linewidth}{0.5pt}\end{center}

\textbf{Sparse sets naturally support many move-only types}

\BookInlineCode{{s}{t}{d}{:}{:}{v}{e}{c}{t}{o}{r}{\textless{}}{T}{\textgreater{}}}{std::vector\textless{}T\textgreater{}} can hold move-only T.

Insertion:

\begin{verbatim}
registry.Emplace<ScriptInstance>(
    entity,
    std::make_unique<VM>());
\end{verbatim}

works if constructors and vector relocation requirements are satisfied.

Swap-removal also uses move semantics.

So a sparse-set ECS can be friendly to move-only components.

\begin{center}\rule{0.5\linewidth}{0.5pt}\end{center}

\textbf{But entity duplication becomes semantically ambiguous}

Suppose entity contains:

\begin{verbatim}
Transform
Health
unique_ptr<ScriptVM>
\end{verbatim}

Now user clicks:

\begin{quote}
Duplicate Entity.
\end{quote}

What should happen to \BookInlineCode{{S}{c}{r}{i}{p}{t}{V}{M}}{ScriptVM}?

You cannot copy the \BookInlineCode{{u}{n}{i}{q}{u}{e}{\_}{p}{t}{r}}{unique\_ptr}.

Possible semantics include:

\[
\text{clone VM state}
\]

or:

\[
\text{create fresh VM from script asset}
\]

or:

\[
\text{do not duplicate runtime-only component}.
\]

This means \textbf{component copyability is not the same as entity duplicability}.

A component metadata system should expose an explicit clone operation.

\begin{center}\rule{0.5\linewidth}{0.5pt}\end{center}

\textbf{Runtime-only components}

A clean classification is:

\begin{verbatim}
persistent authored components
runtime-derived components
runtime-owned external handles
transient frame state
\end{verbatim}

For example:

\begin{verbatim}
LogicComponent:
persistent script asset reference

ScriptRuntimeComponent:
runtime VM instance
\end{verbatim}

Save scenes should serialize the former.

Not the latter.

Duplication may copy the former and reconstruct the latter.

This distinction is already visible in Kairo: \BookInlineCode{{L}{o}{g}{i}{c}{C}{o}{m}{p}{o}{n}{e}{n}{t}}{LogicComponent} stores persistent document identity while VM instances are explicitly process-local.

\begin{center}\rule{0.5\linewidth}{0.5pt}\end{center}

\textbf{Component serialization}

Serialization means converting component state into a durable representation:

\[
Component
\rightarrow
SerializedData.
\]

Deserialization reverses:

\[
SerializedData
\rightarrow
Component.
\]

Do not serialize C++ object memory directly.

For:

\begin{verbatim}
struct Health
{
    float Current;
    float Maximum;
};
\end{verbatim}

you want semantic fields:

\begin{verbatim}
health.current = 80
health.maximum = 100
\end{verbatim}

not bytes at offsets 0 and 4 merely because today's compiler layout happens to look that way.

\begin{center}\rule{0.5\linewidth}{0.5pt}\end{center}

\textbf{Why raw struct dumps are fragile}

Raw dump depends on:

\begin{itemize}
\tightlist
\item
  compiler padding,
\item
  endianness,
\item
  ABI,
\item
  enum representation,
\item
  pointer size,
\item
  alignment,
\item
  implementation changes,
\item
  member order.
\end{itemize}

Suppose:

\begin{verbatim}
struct A
{
    bool Enabled;
    float Value;
};
\end{verbatim}

may be:

\begin{verbatim}
byte 0 bool
bytes 1–3 padding
bytes 4–7 float
\end{verbatim}

A future compiler/configuration can differ.

Durable formats should store logical fields.

\begin{center}\rule{0.5\linewidth}{0.5pt}\end{center}

\textbf{Component serializer interface}

A reflection-independent design:

\begin{verbatim}
struct ComponentSerializer
{
    void (*Serialize)(
        const void* object,
        Writer& writer);

    void (*Deserialize)(
        void* object,
        Reader& reader,
        std::uint32_t version);
};
\end{verbatim}

Type-specific:

\begin{verbatim}
void Serialize(
    const Transform& t,
    Writer& w)
{
    w.Field("position", t.Position);
    w.Field("rotation", t.Rotation);
    w.Field("scale", t.Scale);
}
\end{verbatim}

This is explicit and robust.

\begin{center}\rule{0.5\linewidth}{0.5pt}\end{center}

\textbf{Reflection-driven serialization}

Once reflection metadata exists, generic serialization becomes possible.

If descriptor knows:

\begin{verbatim}
Component: Camera
Properties:
vertical-fov
near-plane
far-plane
projection
...
\end{verbatim}

generic serializer can enumerate them.

That reduces hand-written code.

But beware:

\[
\boxed{
\text{editor-visible property}
\neq
\text{necessarily serialized property}.
}
\]

Some values are derived, transient, editor-only or runtime-only.

Reflection metadata needs flags.

\begin{center}\rule{0.5\linewidth}{0.5pt}\end{center}

\textbf{Kairo's current scene serialization}

Your EngineCore already has durable scene serialization with explicit parsing, validation and versioning.

Current files use headers such as:

\begin{verbatim}
kairo-scene 4
kairo-scene 5
\end{verbatim}

rather than assuming C++ memory layout.

The serializer imposes input limits such as maximum scene size and entity count, validates numeric ranges, validates assets, and uses a temporary file followed by atomic replacement when saving.

That is exactly the kind of defensive scene I/O a production engine should have.

\begin{center}\rule{0.5\linewidth}{0.5pt}\end{center}

\textbf{Runtime ECS entities are explicitly not serialized in Kairo}

Your KairoECS README says the runtime ECS is process-local.

The SceneECS bridge carries:

\begin{verbatim}
struct AuthoredEntityIdentity
{
    kairo::engine::Entity Authored;
};
\end{verbatim}

beside the runtime ECS handle.

This distinction is crucial:

\[
\boxed{
\text{durable authored ID}
\neq
\text{runtime ECS slot/generation}.
}
\]

Excellent design.

\begin{center}\rule{0.5\linewidth}{0.5pt}\end{center}

\textbf{Entity-reference serialization}

Suppose:

\begin{verbatim}
struct Target
{
    Entity TargetEntity;
};
\end{verbatim}

Runtime value might be:

\begin{verbatim}
Index = 142
Generation = 8
\end{verbatim}

You cannot serialize that as the persistent reference.

On next run:

\begin{verbatim}
{142,8}
\end{verbatim}

may not exist at all.

Instead convert:

\[
RuntimeEntity
\rightarrow
PersistentEntityID.
\]

\begin{center}\rule{0.5\linewidth}{0.5pt}\end{center}

\textbf{Save-time remapping}

Suppose:

\begin{verbatim}
Runtime {41,3} ↔ SceneID 1007
Runtime {93,8} ↔ SceneID 1008
\end{verbatim}

Entity 41 targets runtime 93.

Serialized:

\begin{verbatim}
entity 1007
target 1008
\end{verbatim}

Not:

\begin{verbatim}
target index 93 generation 8.
\end{verbatim}

That is the correct durable relation.

\begin{center}\rule{0.5\linewidth}{0.5pt}\end{center}

\textbf{Load-time remapping}

Loading should generally happen in at least two phases.

First create entities:

\begin{verbatim}
persistent ID 1007 -> runtime E1
persistent ID 1008 -> runtime E2
\end{verbatim}

Then resolve references:

\begin{verbatim}
1007.target = lookup(1008)
\end{verbatim}

The key idea:

\[
\boxed{
\text{create identities first,
resolve cross-references afterward}.
}
\]

\begin{center}\rule{0.5\linewidth}{0.5pt}\end{center}

\textbf{Kairo already applies this pattern for parent references}

\BookInlineCode{{S}{c}{e}{n}{e}{E}{C}{S}{M}{i}{r}{r}{o}{r}{:}{:}{E}{x}{t}{r}{a}{c}{t}{(}{)}}{SceneECSMirror::Extract()} uses two passes.

First pass:

\begin{verbatim}
create all ECS entities
store authored→runtime mapping
\end{verbatim}

Second pass:

\begin{verbatim}
resolve hierarchy parent references
\end{verbatim}

This means a child can reference a parent regardless of authored file/entity ordering.

That is exactly the general technique Part 1479 demands.

\begin{center}\rule{0.5\linewidth}{0.5pt}\end{center}

\textbf{Broken references}

What happens when serialized entity A references entity B but B is missing?

Options include:

\begin{verbatim}
hard error
null reference
warning + null
create placeholder
\end{verbatim}

The correct policy depends on the field.

A required parent reference may warrant an error.

An optional target may safely become null.

The serializer/reflection metadata can express:

\begin{verbatim}
required
optional
weak reference
external scene reference
\end{verbatim}

instead of treating every entity reference identically.

\begin{center}\rule{0.5\linewidth}{0.5pt}\end{center}

\textbf{Versioned component schemas}

Suppose Camera v1:

\begin{verbatim}
struct CameraV1
{
    float Fov;
    float Near;
    float Far;
};
\end{verbatim}

Later v2 adds:

\begin{verbatim}
float Exposure;
\end{verbatim}

Old scenes do not contain Exposure.

What should happen?

A schema version allows migration.

\begin{center}\rule{0.5\linewidth}{0.5pt}\end{center}

\textbf{Component versioning model}

Store:

\begin{verbatim}
component Camera
schema 2
...
\end{verbatim}

or version the enclosing scene format.

Then:

\begin{verbatim}
switch (serializedVersion)
{
case 1:
    camera.Fov = read(...);
    camera.Near = read(...);
    camera.Far = read(...);
    camera.Exposure = 0.0f;
    break;

case 2:
    ...
    break;
}
\end{verbatim}

The old format is interpreted intentionally.

\begin{center}\rule{0.5\linewidth}{0.5pt}\end{center}

\textbf{Why a global scene version is sometimes insufficient}

Suppose:

\begin{verbatim}
Scene format = 17
\end{verbatim}

but individual components evolve independently.

Changing Audio should not conceptually require migration logic for Physics.

Component-level schema versions can be cleaner:

\begin{verbatim}
Transform schema 3
Camera schema 5
AudioEmitter schema 2
\end{verbatim}

However, global scene versions are simpler and can still work very well for a tightly controlled engine.

Kairo currently uses whole scene-format versions.

\begin{center}\rule{0.5\linewidth}{0.5pt}\end{center}

\textbf{Migration chains versus direct migration}

Suppose you have:

\begin{verbatim}
v1 → v2 → v3 → v4 → v5.
\end{verbatim}

One model applies every migration sequentially.

Advantages: each migration is small.

Disadvantages: loading very old assets requires all historical migration code forever.

Alternative:

\begin{verbatim}
v1 → current
v2 → current
v3 → current
...
\end{verbatim}

More code duplication but fewer intermediate dependencies.

A common authoring pipeline eventually \textbf{re-saves/cooks old assets into current format}, allowing runtime to support only recent/current versions.

\begin{center}\rule{0.5\linewidth}{0.5pt}\end{center}

\textbf{Kairo's current V5 pattern}

Your V5 serializer is particularly interesting.

It delegates much of the previous format and then extends the result for newer fields such as audio components.

This is a practical evolutionary strategy:

\[
\text{reuse stable previous serializer}
+
\text{append/migrate new semantics}.
\]

It reduces unnecessary rewriting while keeping the version boundary explicit.

\begin{center}\rule{0.5\linewidth}{0.5pt}\end{center}

\textbf{Schema evolution rules}

When designing durable formats, these changes are usually comparatively easy:

\begin{verbatim}
add optional field with default
rename with migration
split one field into two with conversion
change enum with explicit mapping
\end{verbatim}

Harder:

\begin{verbatim}
remove semantic information
change coordinate system
change identity model
merge/split entities
change units
\end{verbatim}

Those require actual migration logic.

Never assume:

\begin{quote}
``We'll just deserialize old bytes into the new struct.''
\end{quote}

That is exactly what schema versioning exists to prevent.

\begin{center}\rule{0.5\linewidth}{0.5pt}\end{center}

\textbf{Reflection-driven component inspection}

Reflection answers:

\begin{quote}
What type is this?
\end{quote}

\begin{quote}
What properties does it expose?
\end{quote}

\begin{quote}
How do I read/write them generically?
\end{quote}

Kairo already has a substantial reflection layer.

For example, it registers descriptors for:

\begin{verbatim}
Transform
Camera
Environment
Light
MeshRenderer
Logic
RigidBody
Collider
\end{verbatim}

with human-readable names, categories, tooltips, numeric ranges, enum options and property read/write functions.

That is significantly beyond toy reflection.

\begin{center}\rule{0.5\linewidth}{0.5pt}\end{center}

\textbf{Property descriptor model}

Conceptually:

\begin{verbatim}
struct PropertyDescriptor
{
    std::string Key;
    std::string DisplayName;

    PropertyValueKind Kind;

    PropertyValue (*Read)(const void*);
    void (*Write)(
        void*,
        const PropertyValue&);

    NumericRange Range;
    PropertyFlags Flags;
};
\end{verbatim}

Now editor does not need:

\begin{verbatim}
if component == Camera
    hardcode every camera field
\end{verbatim}

It can inspect metadata.

\begin{center}\rule{0.5\linewidth}{0.5pt}\end{center}

\textbf{Reflection requires semantic adapters}

Consider:

\begin{verbatim}
Quaternionf Rotation;
\end{verbatim}

Generic bytes are not editor-friendly.

Kairo's reflection explicitly encodes/decodes quaternion/vector property values and normalizes rotation on write.

That is important.

Reflection should expose \textbf{semantic values}, not simply memory addresses.

Likewise asset handles become persistent reference strings/IDs.

\begin{center}\rule{0.5\linewidth}{0.5pt}\end{center}

\textbf{Reflection is not C++ RTTI}

\BookInlineCode{{t}{y}{p}{e}{i}{d}{(}{T}{)}}{typeid(T)} gives some runtime type identity.

It does not give:

\begin{verbatim}
DisplayName
fields
ranges
tooltip
asset target type
editor category
serialization policy
\end{verbatim}

Engine reflection is a domain metadata system.

RTTI and reflection solve different problems.

\begin{center}\rule{0.5\linewidth}{0.5pt}\end{center}

\textbf{Editor property panels}

Given reflected component:

\begin{verbatim}
Camera
\end{verbatim}

editor can generate:

\begin{verbatim}
Projection       [Perspective ▼]
Vertical FOV     [60.0       ]
Near Plane       [0.1        ]
Far Plane        [1000       ]
Exposure         [0.0        ]
Primary          [✓]
\end{verbatim}

from metadata.

No hand-written Camera inspector is strictly required.

\begin{center}\rule{0.5\linewidth}{0.5pt}\end{center}

\textbf{Generic inspector algorithm}

Conceptually:

\begin{verbatim}
for (const PropertyDescriptor& property
     : component.Properties)
{
    PropertyValue value =
        property.Read(componentPointer);

    PropertyValue edited =
        DrawEditorWidget(
            property,
            value);

    if (edited != value)
    {
        property.Write(
            componentPointer,
            edited);
    }
}
\end{verbatim}

The widget depends on property kind.

\begin{center}\rule{0.5\linewidth}{0.5pt}\end{center}

\textbf{Widget mapping}

Examples:

\begin{verbatim}
bool        → checkbox
float       → drag/slider
enum        → combo box
Vec3        → xyz editor
Quaternion  → Euler UI or quaternion editor
Color       → color picker
AssetRef    → asset selector
EntityRef   → scene/entity selector
string      → text box
\end{verbatim}

Reflection metadata provides enough information to choose sensible editors.

\begin{center}\rule{0.5\linewidth}{0.5pt}\end{center}

\textbf{Validation after reflective writes}

Suppose user changes:

\begin{verbatim}
NearPlane = 100
FarPlane  = 10
\end{verbatim}

A generic property write may temporarily make Camera invalid.

You need a policy.

Options:

\begin{verbatim}
validate each property immediately
edit temporary copy then validate/commit
allow transient invalid editor state
clamp through metadata
\end{verbatim}

For complex invariants involving multiple fields, editing a temporary copy then validating is often safest.

\begin{center}\rule{0.5\linewidth}{0.5pt}\end{center}

\textbf{Editor writes should not bypass engine invariants}

Current Kairo components already expose:

\begin{verbatim}
Validate()
\end{verbatim}

for Camera, Light, Collider, etc.

A reflection inspector should not write arbitrary bytes and leave them invalid.

A strong flow:

\begin{verbatim}
read component
    ↓
copy to candidate
    ↓
apply reflected property change
    ↓
candidate.Validate()
    ↓
commit candidate
\end{verbatim}

If validation fails, editor can show an error without corrupting live state.

\begin{center}\rule{0.5\linewidth}{0.5pt}\end{center}

\textbf{Undo/redo for component edits}

Suppose user changes light intensity:

\[
1000\rightarrow5000.
\]

Undo must restore 1000.

Redo must restore 5000.

The simplest command stores:

\begin{verbatim}
struct PropertyEditCommand
{
    PersistentEntityID Entity;
    ComponentSchemaID Component;
    PropertyID Property;

    PropertyValue Before;
    PropertyValue After;
};
\end{verbatim}

\begin{center}\rule{0.5\linewidth}{0.5pt}\end{center}

\textbf{Execute, undo and redo}

\begin{verbatim}
void Execute()
{
    Set(
        Entity,
        Component,
        Property,
        After);
}

void Undo()
{
    Set(
        Entity,
        Component,
        Property,
        Before);
}
\end{verbatim}

Undo stack:

\begin{verbatim}
A
B
C <- current
\end{verbatim}

Undo:

\begin{verbatim}
A
B <- current
C
\end{verbatim}

A new edit after undo usually drops C and begins a new branch.

\begin{center}\rule{0.5\linewidth}{0.5pt}\end{center}

\textbf{Do not store raw component pointers in undo commands}

Bad:

\begin{verbatim}
struct Command
{
    Transform* Target;
};
\end{verbatim}

Component storage may move.

Entity may be destroyed and recreated.

Editor may reload scene.

Use stable authored identity:

\begin{verbatim}
Scene Entity ID
+
Component schema ID
+
Property key
\end{verbatim}

Then resolve on execute/undo.

This is yet another application of:

\[
\text{identity}\neq\text{address}.
\]

\begin{center}\rule{0.5\linewidth}{0.5pt}\end{center}

\textbf{Undo command merging}

Dragging a slider can generate:

\begin{verbatim}
1000
1001
1002
...
5000
\end{verbatim}

You do not want 4000 undo steps.

During one interactive edit transaction:

\begin{verbatim}
BeginEdit
capture Before

many live updates

EndEdit
capture final After
push one command
\end{verbatim}

This is essential for usable editors.

\begin{center}\rule{0.5\linewidth}{0.5pt}\end{center}

\textbf{Structural undo}

Changing a property is easy.

Undoing:

\begin{verbatim}
Delete Entity
\end{verbatim}

requires enough data to recreate:

\begin{itemize}
\tightlist
\item
  persistent ID,
\item
  components,
\item
  hierarchy,
\item
  children policy,
\item
  references,
\item
  asset references.
\end{itemize}

Often the easiest representation is a serialized snapshot of the deleted subtree.

Undo deletion:

\begin{verbatim}
deserialize snapshot
restore IDs/references
\end{verbatim}

This directly connects undo to serialization.

\begin{center}\rule{0.5\linewidth}{0.5pt}\end{center}

\textbf{Prefab/template entities}

A prefab is reusable authored entity structure.

Example:

\begin{verbatim}
EnemySoldier.prefab

Root
├── Transform
├── MeshRenderer
├── RigidBody
├── Collider
├── Health
├── Logic
└── Weapon child
\end{verbatim}

Instantiate it 500 times without manually recreating the structure.

\begin{center}\rule{0.5\linewidth}{0.5pt}\end{center}

\textbf{Prefab is not merely copying an entity}

A useful prefab system needs:

\[
\text{template identity}
+
\text{instance identity}
+
\text{overrides}.
\]

Suppose prefab Health.Maximum changes:

\[
100\rightarrow120.
\]

Instances without overrides should inherit 120.

An instance with override:

\[
Maximum=200
\]

should remain 200.

Thus prefab instances are usually:

\[
\boxed{
Template
+
OverrideSet
}
\]

rather than frozen copies only.

\begin{center}\rule{0.5\linewidth}{0.5pt}\end{center}

\textbf{Prefab instance override model}

Conceptually:

\begin{verbatim}
Prefab:
Health.Maximum = 100
Renderable.Material = SoldierMaterial

Instance:
override Health.Maximum = 200
\end{verbatim}

Resolved value:

\[
Value(instance,p)=
\begin{cases}
Override(p),&p\text{ overridden}\\
Prefab(p),&otherwise
\end{cases}
\]

Nested prefabs add another layer.

This quickly becomes an asset/inheritance system.

\begin{center}\rule{0.5\linewidth}{0.5pt}\end{center}

\textbf{Prefab runtime extraction}

Do not make runtime ECS query template inheritance constantly.

At load/cook/instantiate time:

\begin{verbatim}
Prefab
+ overrides
    ↓
resolved entity data
    ↓
runtime ECS components
\end{verbatim}

Runtime wants concrete component values.

Authoring wants inheritance.

Separate those layers.

\begin{center}\rule{0.5\linewidth}{0.5pt}\end{center}

\textbf{Scene loading into an ECS}

Kairo already has a very clean model for this.

Authoring:

\begin{verbatim}
EngineCore::Scene
\end{verbatim}

Runtime:

\begin{verbatim}
KairoECS::Registry.
\end{verbatim}

Bridge:

\begin{verbatim}
SceneECSMirror::Extract(Scene).
\end{verbatim}

This is much better than serializing ECS dense indices.

\begin{center}\rule{0.5\linewidth}{0.5pt}\end{center}

\textbf{Kairo's two-pass scene extraction}

Current bridge performs:

\textbf{First pass}

For every authored entity:

\begin{verbatim}
create runtime Entity
add AuthoredEntityIdentity
add RuntimeName
add RuntimeTransform
add RuntimeSceneState
record authored→runtime
\end{verbatim}

\textbf{Second pass}

Resolve parent:

\begin{verbatim}
authored parent ID
    ↓
m_byAuthored
    ↓
runtime parent Entity
\end{verbatim}

This is exactly how cross-reference loading should work.

\begin{center}\rule{0.5\linewidth}{0.5pt}\end{center}

\textbf{Why extraction rather than aliasing is valuable}

The bridge test explicitly verifies:

\begin{verbatim}
change authored scene name
\end{verbatim}

does not automatically mutate runtime ECS.

Only:

\begin{verbatim}
RefreshAuthoredState(...)
\end{verbatim}

crosses the boundary.

That gives clear ownership:

\[
\boxed{
Scene = authored source of truth
}
\]

\[
\boxed{
ECS = extracted process-local runtime state.
}
\]

No accidental shared mutable state.

\begin{center}\rule{0.5\linewidth}{0.5pt}\end{center}

\textbf{Scene loading pipeline}

A more complete production pipeline becomes:

\begin{verbatim}
scene file
   ↓
parse
   ↓
validate
   ↓
migrate schema
   ↓
create authored Scene
   ↓
resolve asset references
   ↓
extract runtime ECS
   ↓
instantiate runtime-only components
   ↓
create physics/render/audio bridges
   ↓
BeginPlay
\end{verbatim}

Every arrow is an explicit boundary.

That is good architecture.

\begin{center}\rule{0.5\linewidth}{0.5pt}\end{center}

\textbf{Transactional loading}

Never partially mutate the live world while parsing.

Kairo follows this principle in several places.

Load into:

\begin{verbatim}
candidate Scene
\end{verbatim}

then, if everything succeeds:

\begin{verbatim}
destination = std::move(candidate);
\end{verbatim}

Similarly \BookInlineCode{{R}{u}{n}{t}{i}{m}{e}{W}{o}{r}{l}{d}{:}{:}{C}{o}{m}{m}{i}{t}{(}{)}}{RuntimeWorld::Commit()} applies commands to a copy and only replaces the live scene when all commands succeed.

That provides strong all-or-nothing behavior.

\begin{center}\rule{0.5\linewidth}{0.5pt}\end{center}

\textbf{Scene saving from an ECS}

This is trickier.

If authored Scene remains the source of truth, you generally should not save arbitrary runtime ECS directly.

Runtime may contain:

\begin{verbatim}
PhysicsProxy
ScriptVM
InterpolatedTransform
NavigationAgentRuntime
GPUVisibilityState
\end{verbatim}

that should never become authored scene data.

So saving needs an extraction policy:

\[
Runtime
\rightarrow
\text{persistent authored subset}.
\]

\begin{center}\rule{0.5\linewidth}{0.5pt}\end{center}

\textbf{Runtime-to-authoring sync}

Potential flow:

\begin{verbatim}
for each runtime entity with AuthoredEntityIdentity:
    find authored scene entity

    copy approved editable state:
        Transform
        enabled
        selected gameplay properties

    ignore:
        BodyID
        RenderProxy
        ScriptVM
        transient events
\end{verbatim}

Then serialize the authored Scene.

This whitelist policy is safer than dumping every component.

\begin{center}\rule{0.5\linewidth}{0.5pt}\end{center}

\textbf{Play mode editor semantics}

Editors often distinguish:

\begin{verbatim}
Edit mode
Play mode
\end{verbatim}

During play, runtime transform may change due to physics.

When play stops, should those changes persist?

Possible policy:

\begin{verbatim}
discard all runtime changes
\end{verbatim}

or:

\begin{verbatim}
user explicitly Apply Runtime Changes
\end{verbatim}

or:

\begin{verbatim}
certain fields persist.
\end{verbatim}

That is not purely serialization.

It is product/editor semantics.

But ECS architecture must permit it.

\begin{center}\rule{0.5\linewidth}{0.5pt}\end{center}

\textbf{Entity duplication}

Naive duplication:

\begin{verbatim}
for every component on source:
    copy to destination
\end{verbatim}

sounds easy.

But we immediately encounter categories:

\begin{verbatim}
plain value component
persistent asset reference
entity reference
runtime handle
move-only runtime object
hierarchy relationship
unique identity
\end{verbatim}

They cannot all be copied identically.

\begin{center}\rule{0.5\linewidth}{0.5pt}\end{center}

\textbf{Component clone metadata}

A production component descriptor might contain:

\begin{verbatim}
enum class ClonePolicy
{
    CopyValue,
    Reinitialize,
    Skip,
    Custom
};
\end{verbatim}

Examples:

\begin{verbatim}
Transform       CopyValue
Health          CopyValue
MeshRenderer    CopyValue
PhysicsBodyID   Reinitialize/Skip
ScriptRuntime   Reinitialize
NetworkOwner    Custom
\end{verbatim}

Duplication becomes metadata-driven.

\begin{center}\rule{0.5\linewidth}{0.5pt}\end{center}

\textbf{Entity cloning with remapped references}

Suppose hierarchy:

\begin{verbatim}
A
├── B
└── C
\end{verbatim}

and B contains:

\begin{verbatim}
Target = C.
\end{verbatim}

Duplicate subtree.

Create:

\begin{verbatim}
A'
├── B'
└── C'
\end{verbatim}

What should:

\begin{verbatim}
B'.Target
\end{verbatim}

be?

Usually:

\[
C'
\]

not original C.

This is \textbf{reference remapping}.

\begin{center}\rule{0.5\linewidth}{0.5pt}\end{center}

\textbf{Clone map}

First clone all identities:

\begin{verbatim}
A → A'
B → B'
C → C'
\end{verbatim}

Store:

\begin{verbatim}
std::unordered_map<
    PersistentEntityID,
    PersistentEntityID
> remap;
\end{verbatim}

Then clone components.

For every entity-reference field:

\[
ref'=
\begin{cases}
remap[ref],&ref\text{ inside cloned set}\\
ref,&ref\text{ external and policy preserves it}\\
null/error,&\text{depending policy}
\end{cases}
\]

\begin{center}\rule{0.5\linewidth}{0.5pt}\end{center}

\textbf{Two-pass cloning is the same pattern again}

Notice the repeating architecture:

\textbf{Scene load}

Create all entities → resolve references.

\textbf{Scene→ECS extraction}

Create all runtime entities → resolve parent.

\textbf{Subtree clone}

Create all clones → remap references.

These are all instances of:

\[
\boxed{
\text{identity establishment before relational resolution}.
}
\]

This is a powerful general engine pattern.

\begin{center}\rule{0.5\linewidth}{0.5pt}\end{center}

\textbf{Internal versus external references during clone}

Suppose duplicated Weapon points to shared asset/controller entity outside subtree.

Should it remap?

No, usually.

So clone algorithm must know whether target belongs to cloned closure.

If:

\[
target\in CloneSet,
\]

remap.

Otherwise preserve or clear depending on relationship semantics.

\begin{center}\rule{0.5\linewidth}{0.5pt}\end{center}

\textbf{Parent references during cloning}

Duplicating an entity whose parent is outside the cloned set needs a policy.

For editor ``Duplicate'':

\begin{verbatim}
clone becomes sibling
\end{verbatim}

so both original and duplicate share parent.

For ``Duplicate subtree as prefab'':

root may become detached.

Again, cloning is not one universal byte-copy operation.

\begin{center}\rule{0.5\linewidth}{0.5pt}\end{center}

\textbf{Network entity IDs versus local IDs}

Multiplayer introduces another identity namespace.

Client A runtime entity:

\begin{verbatim}
{Index=81, Generation=3}
\end{verbatim}

Client B representing the same game object:

\begin{verbatim}
{Index=409, Generation=9}.
\end{verbatim}

Therefore ECS handles cannot be network identity.

\begin{center}\rule{0.5\linewidth}{0.5pt}\end{center}

\textbf{Network identity model}

Server assigns:

\begin{verbatim}
NetworkEntityID = 761283.
\end{verbatim}

Each machine maintains:

\[
NetworkID
\leftrightarrow
LocalEntity.
\]

For client:

\begin{verbatim}
761283 -> {409,9}
\end{verbatim}

For server:

\begin{verbatim}
761283 -> {81,3}.
\end{verbatim}

Replication packets mention:

\begin{verbatim}
761283
\end{verbatim}

never process-local ECS slot IDs.

\begin{center}\rule{0.5\linewidth}{0.5pt}\end{center}

\textbf{Network ID lifetime}

Network IDs themselves need lifetime rules.

If IDs are reused rapidly, delayed packets can create another ABA problem.

Solutions include:

\begin{itemize}
\tightlist
\item
  monotonically increasing large IDs,
\item
  epoch/generation,
\item
  connection-local spawn sequence,
\item
  UUID-like IDs,
\item
  destroy acknowledgement before reuse.
\end{itemize}

The same identity lesson repeats at another scale.

\begin{center}\rule{0.5\linewidth}{0.5pt}\end{center}

\textbf{Mapping creation order}

Suppose packet says:

\begin{verbatim}
Entity 500 references Entity 700.
\end{verbatim}

but 700's spawn packet has not arrived.

You need:

\begin{verbatim}
unresolved reference
\end{verbatim}

or deferred packet application.

Network replication is effectively distributed entity-reference deserialization.

So much of Part 1479's logic reappears in networking.

\begin{center}\rule{0.5\linewidth}{0.5pt}\end{center}

\textbf{Deterministic entity iteration}

Suppose physics/gameplay depends on entity update order.

If iteration order changes between runs, floating-point accumulation and event order can diverge.

Sparse-set order is affected by:

\begin{verbatim}
insertion
deletion
swap-remove.
\end{verbatim}

Therefore:

\[
\boxed{
\text{dense order is not inherently stable semantic order}.
}
\]

\begin{center}\rule{0.5\linewidth}{0.5pt}\end{center}

\textbf{Swap-remove destroys ordering}

Start:

\begin{verbatim}
A B C D
\end{verbatim}

remove B:

\begin{verbatim}
A D C.
\end{verbatim}

So D now executes before C.

If gameplay behavior depended on:

\begin{verbatim}
A B C D
\end{verbatim}

it is broken.

The ECS is correct.

The gameplay assumption was wrong.

\begin{center}\rule{0.5\linewidth}{0.5pt}\end{center}

\textbf{Determinism requires an explicit contract}

Options include:

\begin{verbatim}
sort by stable entity ID before order-sensitive work
\end{verbatim}

or:

\begin{verbatim}
maintain ordered storage
\end{verbatim}

or:

\begin{verbatim}
make system semantics order-independent
\end{verbatim}

or:

\begin{verbatim}
define deterministic swap-remove lifecycle.
\end{verbatim}

The best solution is often to make most systems order-insensitive, then impose order only where necessary.

\begin{center}\rule{0.5\linewidth}{0.5pt}\end{center}

\textbf{Kairo authored Scene already gives deterministic iteration}

\BookInlineCode{{S}{c}{e}{n}{e}{:}{:}{E}{n}{t}{i}{t}{i}{e}{s}{(}{)}}{Scene::Entities()} gathers IDs then sorts them in ascending stable scene-ID order.

That is excellent for:

\begin{itemize}
\tightlist
\item
  serialization,
\item
  deterministic editor behavior,
\item
  stable diff output.
\end{itemize}

KairoECS dense pool order is instead runtime-storage order.

Those differing goals justify differing iteration contracts.

\begin{center}\rule{0.5\linewidth}{0.5pt}\end{center}

\textbf{Floating-point determinism}

Even identical entity order does not automatically provide bit-identical simulation across:

\begin{verbatim}
x86
ARM
different compilers
different SIMD modes
different thread schedules.
\end{verbatim}

Floating-point reductions are non-associative:

\[
(a+b)+c
\neq
a+(b+c)
\]

in general finite precision.

Parallel reduction order can change results.

For strict lockstep determinism, you must control much more than ECS order.

Do not promise ``deterministic ECS'' and imply ``deterministic entire engine.''

\begin{center}\rule{0.5\linewidth}{0.5pt}\end{center}

\begin{center}\rule{0.5\linewidth}{0.5pt}\end{center}

\hyperref[volume-reading-map]{Reading map} · \hyperref[05-engine-i-ch016-archetypes-chunks-relationships-and-transforms]{Previous chapter} · \hyperref[05-engine-i-ch018-ecs-contracts-diagnostics-and-gpu-handoff]{Next chapter}
